\section{Relative traces and the span of the periods}
\label{sec:trace}

Fix the family selected in Section~\ref{sec:assembly}, a scale \(x\)
beyond its thresholds, and the record at one of the tied translates.
Write \(W\) for its finite set of words. The two full Gram estimates and
the multiplicity estimate are the inputs to this section's conditional
argument. We use them to control the span of the periods after restricting
each column space and projecting onto the common row space. The
orbital calculations give context for the Gram estimates and the
same-word estimate also explains the construction of the nonzero
bottom reference pairing in Section~\ref{sec:periods}.

\begin{figure}[!htbp]
\centering
\begin{tikzpicture}[x=1mm,y=1mm]
  \node[argument block,fill=TraceTint] (diagonal) at (0,0)
    {\hyperlink{proof.full-diagonal-estimate}{\shortstack{Diagonal Gram bound\\\(\|K_{ww}-I\|_{\mathrm{HS}}^2\le e_0 d_w\)}}};
  \node[argument block] (anchor) at (55,0)
    {\hyperlink{proof.nonzero-anchor}{\shortstack{Reference pairing\\Optional proof route}}};
  \node[argument block] (periods) at (110,0)
    {\hyperlink{proof.rational-period-map}{\shortstack{Period maps \(B_{K,w}\)\\with a common target}}};
  \node[argument block,fill=TraceTint] (mixed) at (0,-20)
    {\hyperlink{proof.mixed-word-estimate}{\shortstack{Off-diagonal bounds\\for \(K_{wv}\), \(w\ne v\)}}};
  \node[argument block,fill=TraceTint] (compression) at (55,-20)
    {\hyperlink{proof.compressed-frame-moments}{\shortstack{Trace moments after\\column restriction}}};
  \node[argument block,fill=TraceTint] (rational) at (110,-20)
    {\hyperlink{proof.rational-corank}{\shortstack{Codimension bound\\in the row dual \(C_K\)}}};
  \draw[argument arrow,densely dashed] (diagonal.east) -- (anchor.west);
  \draw[argument arrow,densely dashed] (anchor.east) -- (periods.west);
  \draw[argument arrow] (diagonal.south) -- ++(0,-4) -| (compression.north);
  \draw[argument arrow] (mixed.east) -- (compression.west);
  \draw[argument arrow] (compression.east) -- (rational.west);
  \draw[argument arrow] (periods.south) -- (rational.north);
\end{tikzpicture}
\caption{The dashed route gives the reference pairing for an optional
  proof of the granted simultaneous period comparison. There the diagonal
  and ramified rank bounds give \(\dim\ker\mathcal P_w\le e_0d_w<r_w\),
  and full row evaluation pairs with a nonzero image of a ramified
  column. To descend the granted comparison over \(K\), use the compatible
  base pairing and the surjective projected row map over \(K\).
  With the word and dimension estimates, both Gram bounds control the
  trace moments for isometric column inclusions. Using \(J_w=\jmath_w U_w\)
  matches those moments to the period maps. The common projection and
  the distinguished-row dimension bound give the corank over \(K\).}
\label{fig:trace-handoffs}
\end{figure}


\subsection{Unfolding the full Gram operators}
\leavevmode\teamanchor{proof.full-gram-normalization}%
For each \(w\in W\), use the Poincar\'e-series operator
\(\mathcal P_w:H_w^{\mathrm{full}}\to\mathcal H\) in
\eqref{eq:period-gram}. The ambient level and archimedean row type fix
the same finite-dimensional Petersson space \(\mathcal H\) for all
words. Let \(d_w=\dim H_w^{\mathrm{full}}\), \(D=\dim\mathcal H\), and
\(R_{\mathrm{full}}=\sum_w d_w\). The coefficient representation
\(\pi_w\) is the tensor product of the residual compact \(U(k)\)-types
at the real places. Write \(v_N(w)\) for the volume of the unipotent
fundamental domain multiplied by the ambient level volume, and
\(v_H(w)\) for the volume of the real Bessel stabilizer. Then
\(c_0(w)=v_N(w)v_H(w)/\dim\pi_w\), with the measures of
\eqref{eq:period-gram}. Set \(L_w=c_0(w)^{-1/2}\mathcal P_w\) and
\(K_{wv}=L_w^*L_v:H_v^{\mathrm{full}}\to H_w^{\mathrm{full}}\).
The inner product is linear in its second argument.

The orbital formulas in the next paragraphs describe the operators
in the two Gram inputs. Their quantitative estimates are used in the
forms stated at \hyperlink{candidate.diagonal}{the same-word estimate}
and \hyperlink{candidate.trace}{the mixed estimate}.

\leavevmode\teamanchor{proof.gram-unfolding}%
Put \(G_0=U(Y_0)\), where \(Y_0\) has rank \(r=k+m-1\) over the CM
field \(E\), and write \(\mathcal S_w\) for the Bessel subgroup attached to
the flag and its distinguished vector for \(w\). For \(w,v\in W\),
\(\phi\in H_w^{\mathrm{full}}\) and \(\psi\in H_v^{\mathrm{full}}\),
the corresponding Poincar\'e series periodize the seeds
\(\Phi_{w,\phi}\) and \(\Phi_{v,\psi}\), which are invariant under their
respective left rational Bessel subgroup actions, over
\(\mathcal S_w(F)\backslash G_0(F)\) and
\(\mathcal S_v(F)\backslash G_0(F)\), respectively.
Let \(\Omega_{wv}\) be representatives of the rational double cosets
\(\mathcal S_w(F)\backslash G_0(F)/\mathcal S_v(F)\), and set
\(\mathcal S_\gamma(F)=\mathcal S_w(F)\cap\gamma\mathcal S_v(F)\gamma^{-1}\).
Put
\[
 I^+_{wv}(\phi,\psi)=\sum_{\gamma\in\Omega_{wv}}
   \int_{\mathcal S_\gamma(F)\backslash G_0(\mathbf A_F)}
    |\Phi_{w,\phi}(g)|\,|\Phi_{v,\psi}(\gamma^{-1}g)|\,dg,\qquad
 P^+_{w,\phi}(g)=\sum_{\eta\in\mathcal S_w(F)\backslash G_0(F)}
    |\Phi_{w,\phi}(\eta g)|,
\]
with values in \([0,\infty]\). The stabilizer reindexing described
below uses only the left invariance of \(|\Phi_{w,\phi}|\) under
\(\mathcal S_w(F)\) and of \(|\Phi_{v,\psi}|\) under \(\mathcal S_v(F)\).
For nonnegative functions Tonelli's theorem allows it without any
convergence hypothesis, so
\(I^+_{wv}(\phi,\psi)=\int_{G_0(F)\backslash G_0(\mathbf A_F)}
P^+_{w,\phi}P^+_{v,\psi}\,dg\).
By the Cauchy--Schwarz inequality,
\(I^+_{wv}(\phi,\psi)^2\le I^+_{ww}(\phi,\phi)\,I^+_{vv}(\psi,\psi)\).
Absolute convergence for all pairs therefore reduces to the diagonal
pairs, that is, to square integrability of the absolute Poincar\'e
series \(P^+_{w,\phi}\). The edition takes this diagonal finiteness from
the release: at the chosen level ai agent proves
\releasecite{Formalism/LocalToGlobal/SiegelLatticeTheory/FlagStabilizerTensors/LinkSFinal.lean}{97}{absolute convergence}
using the level conditions, neatness and real Whittaker decay.
The release states this finiteness for every pair of words; only its
diagonal instance is used here.
Hence \(I^+_{wv}(\phi,\psi)<\infty\) for all \(w,v,\phi,\psi\). In
particular both Poincar\'e series converge absolutely almost
everywhere, and Fubini's theorem justifies every exchange of sum and
integral in the reindexing. The double-coset unfolding therefore
gives the normalized identity
\[
 \langle\phi,K_{wv}\psi\rangle
   =\frac{1}{\sqrt{c_0(w)c_0(v)}}\sum_{\gamma\in\Omega_{wv}}
     \int_{\mathcal S_\gamma(F)\backslash G_0(\mathbf A_F)}
      \overline{\Phi_{w,\phi}(g)}\Phi_{v,\psi}(\gamma^{-1}g)\,dg .
\]
All quotients carry the measures induced by the adelic Haar measure
in the Petersson product. After unfolding the \(w\) periodization,
the \(v\) sum is indexed by \(\mathcal S_v(F)\backslash G_0(F)\).
The cosets in the double coset of \(\gamma^{-1}\) are represented
by \(\gamma^{-1}a\) for
\(a\in\mathcal S_\gamma(F)\backslash\mathcal S_w(F)\). Integrating
their translates joins the corresponding fundamental domains.
Getz--Hahn use the same stabilizer reindexing for the geometric
side of their relative trace formula with compactly supported tests.
\cite[\S4C, (4C.8) and its unfolding, p.~114]{tr3getzhahn}
ai agent
\releasecite{Formalism/LocalToGlobal/SiegelLatticeTheory/FlagStabilizerTensors/LinkSFinal.lean}{92}{identifies the unfolded coefficients}
using compatible flag and lattice representatives. Separating the
lattices and the finite level conditions gives \(K_{wv}=\sum_t A_t\). In each
coefficient block allowed by the level, \(A_t\) is the archimedean
orbital operator multiplied by the finite character, the level volume
and \((c_0(w)c_0(v))^{-1/2}\). The archimedean operator integrates
\(W_w(g)\otimes W_v(\gamma^{-1}g)^*\), where \(W_w,W_v\) are the
coefficient vectors of the Poincar\'e seeds.

\leavevmode\teamanchor{proof.gram-orbit-expansion}%
Each decorated flag comprises an ordered isotropic frame of length
\(r_S=m/2-1\) and a vector orthogonal to that frame of norm
\(-T_w\) or \(-T_v\). Here \(m\) is even and \(m\ge4\) on the
chosen large-rank family. Set \(a=r_S+1\), and write
\(Z\) for the \(a\)-by-\(a\) matrix of cross pairings of the augmented
frames. With \(e\) the last coordinate vector, their joint Gram matrix
and its determinant are
\[
  M=\begin{pmatrix}-T_w ee^*&Z\\ Z^*&-T_v ee^*\end{pmatrix},
  \qquad
  \det M=(-1)^a G,\qquad
  G=|\det Z|^2-T_wT_v|\det Z_{a-1}|^2 ,
\]
where \(Z_{a-1}\) is the leading \((a-1)\)-by-\((a-1)\) submatrix
and \(|z|^2=z\bar z\) is the CM norm. The determinant follows by swapping
the block columns, taking the Schur complement of \(Z\) and using the
cofactor \(e^*\operatorname{adj}(Z)e\). Polynomiality covers singular
\(Z\). When \(G\ne0\) the augmented
frames span a nondegenerate space of dimension \(2a=m\), since \(m\)
is even. Its orthogonal complement has dimension \(k-1\).

Separate the contribution of the identity double coset, the other
\(G=0\) terms, and the generic and nongeneric terms with \(G\ne0\).
Write \(g_0,g_1\) for the two designated auxiliary rational primes,
which split completely in \(E\). At places above these primes the
finite level uses specific Whittaker characters. A nondegenerate pair
is generic if its lattice in the joint span is generated locally by
the two augmented frames, with the adjusted second frame at the
places above the auxiliary primes. At the identity, the level
calculation for \(w=v\) matches columns and gives \(v_N(w)\). The
average of the unit coefficient vector's rank-one projector over
the real stabilizer commutes with
\(\pi_w\), and its trace is \(v_H(w)\). Schur's lemma gives
\((v_H(w)/\dim\pi_w)I\), so the identity contributes \(I\) to
\(K_{ww}\). For distinct words the calculation at the identity gives
zero by averaging between inequivalent coefficient types.

The unipotent average on the stabilizer quotient vanishes when the
characters disagree on the intersection after transport by \(\gamma\).
Let \(W\) be the joint span of the two augmented frames of a compatible
\(G=0\) orbit and let \(R\) be its radical. The pointwise stabilizer
of \(W\) has reductive quotient \(U(r-\dim W-\dim R)\), since
\(r-\dim W-\dim R\) is the rank of the form on \(W^\perp\). Since \(W\)
contains each augmented frame \(F\), we have \(W^\perp\subseteq
F^\perp\), and this rank is at most the rank \(k\) of the form on
\(F^\perp\). Hence \(\dim W+\dim R\ge m-1\). In the case of equality the
image of \(W^\perp\) in \(F^\perp/\operatorname{rad}F^\perp\) has rank
\(k\), so it is the whole space and the kernel is
\(\operatorname{rad}W^\perp\). The induced map from
\(W^\perp/\operatorname{rad}W^\perp\) is an isometry of
\(k\)-dimensional definite spaces at each real place. The pointwise
stabilizer of \(W\) then maps onto the full residual \(U(k)\) of each
flag, up to fixed unitary conjugation. For such an orbit
Schur's lemma forces the resulting intertwiner to vanish for \(w\ne v\)
because the chosen highest weights differ at every real factor. The next
paragraph shows that equality \(\dim W+\dim R=m-1\) holds for every
compatible orbit with \(G=0\), so this argument covers all of them.

\leavevmode\teamanchor{proof.degenerate-normal-form}%
\paragraph{Compatible orbits with \(G=0\).}
Put \(s=r_S\), so \(m=2s+2\) and \(a=s+1\), and write the two augmented
frames, after transport by \(\gamma\), as \((e_1,\dots,e_s,z)\) and
\((e'_1,\dots,e'_s,z')\) with \(\langle z,z\rangle=-T_w\) and
\(\langle z',z'\rangle=-T_v\). The Hermitian form on \(Y_0\) is taken
linear in the first variable. Put \(X_b=\langle e_1,\dots,e_b\rangle\),
\(X=X_s\), \(X^+=X+Ez\) and \(e_{s+1}=z\), and define \(X'_b\), \(X'\),
\(X'^+\), \(e'_{s+1}=z'\) in the same way. The unipotent group \(N\) of
the first flag consists of the isometries that preserve each \(X_b\),
act trivially on each \(X_b/X_{b-1}\) and act trivially on
\(X^\perp/X\). For \(n\in N\) and \(1\le b\le s\) let \(\lambda_b(n)\)
be the \(e_b\)-coordinate of \(ne_{b+1}-e_{b+1}\in X\). The generic
character has the form \(\psi_N(n)=\psi_E(\sum_b\lambda_b(n))\) for a
nontrivial character \(\psi_E\) of \(\mathbf A_E/E\), as in the release's
definition of compatibility. Only the nontriviality of \(\psi_E\) is
used, and nonzero constant multiples of the \(\lambda_b\) would not
change the argument. Define \(N'\) and \(\psi_{N'}\) in the same way.
Compatibility means \(\psi_N=\psi_{N'}\) on
\((N\cap N')(\mathbf A_F)\). The ambient space has dimension
\(r=k+2s+1\) and negative index \(m/2=s+1\), so signature
\((k+s,s+1)\), at each real place. Choose an isotropic \(X^*\) in
duality with \(X\) and orthogonal to \(z\). Then \(X+X^*\) has
signature \((s,s)\), and the orthogonal complement of \(X+X^*+Ez\) has
signature \((k,0)\). Hence the form on \((X^+)^\perp\) is positive
semidefinite at each real place, with radical \(X\) and positive
definite \(k\)-dimensional quotient. Consequently \((\mathrm P)\) every
isotropic rational vector of \((X^+)^\perp\) lies in \(X\), and no
rational vector of \((X^+)^\perp\) has norm \(-T\) with \(T\) totally
positive. The same holds for the second flag. For \(0\ne x\in X\) the
level is the least \(\ell\) with \(x\in X_\ell\), and the leading
coefficient is its \(e_\ell\)-coordinate. Levels and leading
coefficients for the second flag are defined in the same way.

\emph{Step 1.} Let \(0\ne x\in X\cap X'\) have levels \(\ell,\ell'\) and
leading coefficients \(\alpha,\alpha'\). Then
\(\alpha e_{\ell+1}-\alpha'e'_{\ell'+1}\in X_\ell+X'_{\ell'}\). To see
this, let \(y\) be an adelic vector orthogonal to \(X_\ell+X'_{\ell'}\),
choose \(\eta\in\mathbf A_E\) with \(\eta+\bar\eta=-\langle y,y\rangle\),
and put \(g(u)=u+\langle u,y\rangle x-\langle u,x\rangle y
+\eta\langle u,x\rangle x\). Since \(x\) is isotropic and
\(\langle x,y\rangle=0\), expansion gives
\(\langle gu,gu'\rangle=\langle u,u'\rangle
+\langle u,x\rangle\overline{\langle u',x\rangle}
(\eta+\bar\eta+\langle y,y\rangle)=\langle u,u'\rangle\), so \(g\) is
an isometry. For \(b\le s\) we have \(\langle e_b,x\rangle=0\), so
\(ge_b=e_b+\langle e_b,y\rangle x\). This is \(e_b\) for \(b\le \ell\)
and lies in \(e_b+X_\ell\subseteq e_b+X_{b-1}\) for \(b>\ell\). For
\(u\in X^\perp\), \(gu=u+\langle u,y\rangle x\in u+X\). Hence \(g\in N\),
and likewise \(g\in N'\). Since \(gz-z=\langle z,y\rangle x\) and
\(ge_{b+1}-e_{b+1}=\langle e_{b+1},y\rangle x\), with
\(\langle e_{b+1},y\rangle=0\) for \(b<\ell\) and the \(e_b\)-coordinate
of \(x\) zero for \(b>\ell\), we get
\(\psi_N(g)=\psi_E(\langle\alpha e_{\ell+1},y\rangle)\) and
\(\psi_{N'}(g)=\psi_E(\langle\alpha'e'_{\ell'+1},y\rangle)\). Put
\(u_0=\alpha e_{\ell+1}-\alpha'e'_{\ell'+1}\). Compatibility gives
\(\psi_E(\langle u_0,y\rangle)=1\) for all such \(y\). If a rational
\(y\perp X_\ell+X'_{\ell'}\) had \(\langle u_0,y\rangle\ne0\), replacing
\(y\) by \(ty\) with \(t\in\mathbf A_E\) would make \(\psi_E\) trivial.
Hence \(u_0\perp(X_\ell+X'_{\ell'})^\perp\), and
\(u_0\in X_\ell+X'_{\ell'}\) because the ambient form is nondegenerate.

\emph{Step 2.} First, \(z\notin X+X'\): if \(z=x_1+x_1'\) with
\(x_1\in X\) and \(x'_1\in X'\), then \(x'_1=z-x_1\perp X\) and
\(-T_w=\langle z,z\rangle=\langle x'_1,x'_1\rangle=0\). Likewise
\(z'\notin X+X'\). In Step 1 suppose \(\ell,\ell'<s\), and write
\(u_0=p+p'\) with \(p\in X_\ell\) and \(p'\in X'_{\ell'}\). Then
\(\alpha e_{\ell+1}-p=\alpha'e'_{\ell'+1}+p'\) lies in \(X\cap X'\) and
has levels \(\ell+1,\ell'+1\). If \(\ell=s>\ell'\), Step 1 puts
\(\alpha z\) in \(X+X'\), which is impossible, and \(\ell'=s>\ell\) is
excluded in the same way. Iterating from any nonzero \(x\in X\cap X'\)
therefore gives \(\ell=\ell'\) for every such \(x\), and at the top
level gives \(z'-\theta z\in X+X'\) for some \(\theta\in E^\times\)
whenever \(X\cap X'\ne0\). Let \(j+1\) be the least level of a nonzero
element of \(X\cap X'\), with \(j=s\) if \(X\cap X'=0\). The iteration
produces elements of every level \(j+1,\dots,s\), which are
independent, while \(X\cap X'\cap X_j=0\). Hence
\(\dim(X\cap X')=s-j\), \(X=X_j\oplus(X\cap X')\) and, since the levels
agree, \(X'=X'_j\oplus(X\cap X')\).

\emph{Step 3, \(X\cap X'\ne0\).} Let \(Z_j\) be the leading
\(j\)-by-\(j\) submatrix of \(Z\). The pairing of \(X_j\) with \(X'_j\)
is perfect, so \(\det Z_j\ne0\). Indeed, if \(x\in X_j\) is orthogonal
to \(X'_j\), then \(x\perp X\cap X'\), hence \(x\perp X'\), and
\(x\perp z'\) because \(z'-\theta z\in X+X'\). By \((\mathrm P)\) for
the second flag, \(x\in X'\), so \(x\in X\cap X'\cap X_j=0\). Thus
\(Q=X_j+X'_j\) is nondegenerate of dimension \(2j\). Choose \(h\in X_j\)
with \(\tilde z=z-h\perp X'_j\). Then \(\tilde z\perp Q\),
\(\tilde z\perp X\cap X'\) and \(\langle\tilde z,\tilde z\rangle=-T_w\).
Since \(z'\in\theta z+X+X'\), we get
\(W=Q\oplus(X\cap X')\oplus E\tilde z\), orthogonally, with
\(X\cap X'\) isotropic. So \(R=X\cap X'\), \(\dim W=s+j+1\) and
\(\dim R=s-j\).

\emph{Step 3, \(X\cap X'=0\).} Here compatibility is not needed. If
\(\det Z=0\), some nonzero \(u=x+\beta z\) with \(x\in X\) is
orthogonal to \(X'^+\). If \(\beta=0\), \((\mathrm P)\) gives \(x\in
X'\), so \(x=0\). Otherwise we may take \(\beta=1\), and then
\(\langle u,u\rangle=-T_w\), contrary to \((\mathrm P)\). So
\(\det Z\ne0\), and \(G=0\) forces \(\det Z_s\ne0\), where
\(Z_s=Z_{a-1}\). Then \(X\) and \(X'\) are in perfect pairing, and
\(Q=X\oplus X'\) is nondegenerate. Let \(\tilde z=z-h\) and
\(\tilde z'=z'-h'\), with \(h\in X\) and \(h'\in X'\), be the components
orthogonal to \(Q\). Their norms are \(-T_w\) and \(-T_v\). The Gram
determinant of the augmented frames is \(\pm|\det Z_s|^2\) times that
of \((\tilde z,\tilde z')\). So \(G=0\) makes the latter singular, and
some \(y=\tilde z'-\mu\tilde z\) is orthogonal to \(\tilde z\) and
\(\tilde z'\). The coefficient of \(\tilde z'\) cannot vanish, because
\(\tilde z\) is not isotropic. Then \(y\) is isotropic, orthogonal to
\(Q\) and to \(z\), so \(y\in X\subseteq Q\) by \((\mathrm P)\), and
therefore \(y=0\). Hence \(W=X\oplus X'\oplus E\tilde z\) is
nondegenerate, \(R=0\) and \(\dim W=2s+1\).

In both cases \(\dim W+\dim R=2s+1=m-1\), and the Schur argument above
gives the vanishing for \(w\ne v\).

The mixed block
 consists of the generic and nongeneric sums with
\(G\ne0\) in the same weights and measures.

\subsection{The single-word estimate}

Put \(f=[F:\mathbf Q]\) and \(b=\exp(x^\varepsilon)\). The reference
dimension is
\[
 D_\star=\frac{\mathfrak d V_G}{v^\dagger}>0,
\]
where \(\mathfrak d\) is the product of the target formal degrees,
\(V_G\) is the ambient adelic volume and \(v^\dagger\) is the deep
ambient level volume. These are the measures of the Bessel realization.
For the common factor \(\Gamma=\exp(C_\Gamma m^5 f\log b)\), with
\(C_\Gamma>0\) fixed on the family, the column-size output of the
selection in Section~\ref{sec:assembly} is
\begin{equation}
 \label{eq:trace-balance}
 C_0^{-1}\le \frac{d_w\Gamma^2}{D_\star}\le C_0
 \qquad(w\in W), \qquad C_0\ge1.
\end{equation}
The selected constant is \(C_0=3\). The balanced dimensions here count
the full level-zero columns before cutting by the word-prime
characters. Since \(D_\star>0\), the lower bound makes every \(d_w\)
positive. The selected rank inequalities give \(m^4\le f\le k\le x^5\)
and \(k\ge 8\log\Gamma\) at all sufficiently large scales.

The level index explains the balance. Let \(d_w^{(0)}\) be the compact
column dimension before marking a nonzero covector at each place
above \(T_w\), with the two auxiliary levels and the word's real
type retained. Let \(K_0\) be this compact level and \(K_0'\subset K_0\)
the subgroup of elements whose reduction at each place above \(T_w\)
fixes the marked nonzero covector. Reduction onto
\(\mathrm{GL}_k(\mathbf F_{T_w})\) is surjective, and this group acts
transitively on the \(T_w^k-1\) nonzero covectors. By the
orbit--stabilizer theorem, \([K_0:K_0']=(T_w^k-1)^f\).
For a class \(G(F)gK_0\), the classes of level \(K_0'\) above it are
in bijection with the orbits of \(g^{-1}\Gamma_gg\) on \(K_0/K_0'\),
where \(\Gamma_g=G(F)\cap gK_0g^{-1}\). This group is finite by
definiteness.
At a place above the auxiliary prime \(g_i\), the level consists of
elements with upper unitriangular reduction. The preimage of the
unitriangular group is pro-\(g_i\), because the principal congruence
kernel is pro-\(g_i\), and so is each conjugate of it. Hence every
element of \(\Gamma_g\) has order a power of \(g_0\) and a power of
\(g_1\), so \(\Gamma_g=1\). The same holds for the stabilizers at level
\(K_0'\). Thus every class of level \(K_0\) has exactly
\((T_w^k-1)^f\) classes above it, and every class at either level
carries the full residual fibre. Counting fibres gives
\[
 d_w=(T_w^k-1)^f\,d_w^{(0)} .
\]
For comparison, the release records this index as
\releasecite{Formalism/SharpRestriction/GlobalSchurBounds/StandardLattice.lean}{446}{The level index and compact rank}.
Write
\[
 Y_w=\left(\frac{D_\star}{d_w^{(0)}\Gamma^2}\right)^{1/(fk)}.
\]
The word-prime interval is \(Y_w\le T_w\le Y_w(1+1/(fk))\). It is
the auxiliary-level choice of Section~\ref{sec:integral}, which the
edition takes from the release there. In that choice
\(d_w^*\) is the coefficient of \((T_w^k-1)^f\) in the full column
dimension, so \(d_w^*=d_w^{(0)}\) by the index formula, and the chosen
interval \(1\le T_w/(X_w\Gamma^{-2/(fk)})\le1+1/(fk)\) reads
\(Y_w\le T_w\le Y_w(1+1/(fk))\) because
\(X_w\Gamma^{-2/(fk)}=Y_w\). The release records the same interval in
the chosen record at
\releasecite{Formalism/LocalRingCardinality/RhoDetPrimeEqMuTildeMulChiPow.lean}{153}{The word-prime interval}.
Thus the exact ratio \(d_w\Gamma^2/D_\star\) is
\((T_w^k-1)^f/Y_w^{fk}=(T_w/Y_w)^{fk}(1-T_w^{-k})^f\). Its upper bound
is \((1+1/(fk))^{fk}\le\exp(1)<3\). By Bernoulli's inequality its lower
bound is \(1-fT_w^{-k}\ge1-x^5b^{-k}\ge1/2>1/3\) for large \(x\), since
\(T_w\ge b\) and \(f\le x^5\). Given the interval, this proves
\eqref{eq:trace-balance} with \(C_0=3\).
The count uses the same \(D_\star\) as the row projector below.

\leavevmode\teamanchor{proof.full-diagonal-estimate}%
With the contribution \(I\) of the identity removed, there are two
sums in \(K_{ww}-I\): the sum \(K_{\rm deg}\) over the nonidentity
orbits with \(G=0\), and the sum \(K_{\rm nd}=K_{\rm gen}+K_{\rm ng}\)
over the generic and nongeneric orbits with \(G\ne0\). We show
\(\|K_{\rm deg}\|_{\rm HS}^2\le d_w\exp(-kf\log k)\) and
\(\|K_{\rm nd}\|_{\rm HS}^2\le\Gamma d_w^2/D_\star\). Throughout,
\(\log P_g\le6b^{-7}\), as in the
\hyperlink{proof.mixed-word-estimate}{mixed-word estimate}, so
\(P_g\le2\).

\paragraph{The degenerate sum.}
\teamanchor{proof.degenerate-diagonal-sum}%
Group the nonidentity \(G=0\) terms into rows as in the
\hyperlink{proof.quadratic-row-contraction}{quadratic row contraction},
with \(w=v\). The
\hyperlink{proof.squared-norm-identity}{derivation of the squared norm}
applies verbatim to this family \(\mathcal S\). It uses the character
factorization, the constant unipotent factor and the free column
decompositions, which concern single terms, their lattice pairs and
their column representatives and not the family. It also uses the left
and right Gram invariance, whose cited release statements are made
for each of the generic, nongeneric and nonidentity degenerate
families, and the absolute summability, which covers all orbit terms.
Hence
\(\|K_{\rm deg}\|_{\rm HS}^2=\operatorname{Re}\sum_{i,j}\mathcal G_{ij}
\langle V_i,V_j\rangle\), with \(V_i\) formed from the degenerate terms.
The Gelfand--Graev bound
\(0\le\mathcal G\le P_g\Gamma_1\operatorname{Id}\) proved there gives,
by the same coordinate argument,
\[
 \|K_{\rm deg}\|_{\rm HS}^2
 \le P_g\Gamma_1\sum_{\substack{t,t'\ {\rm deg}\\ r(t)=r(t')}}
     \|v_t\|\,\|v_{t'}\|.
\]
The right side satisfies the following typed amplitude bound, which is
assembled below from release inputs. There is a constant \(C_A\ge1\),
fixed on the family, such that
\[
 \Gamma_1\sum_{\substack{t,t'\ {\rm deg}\\ r(t)=r(t')}}
     \|v_t\|\,\|v_{t'}\|
 \le d_w\Bigl(\sum_{j=1}^{r_S}\alpha_j\Bigr)^2,\qquad
 \alpha_j=\exp(C_Am^6f\log b)\Bigl(\frac{C_A}{k}\Bigr)^{j(k-1-j)f}.
\]
The numerics after this bound are given at the end of the assembly.

\leavevmode\teamanchor{proof.degenerate-amplitude-assembly}%
\paragraph{Assembly of the typed amplitude bound.}
Call a degenerate term relevant if its orbit is compatible in the sense
of the \hyperlink{proof.degenerate-normal-form}{normal form}. The
assembly uses six inputs.

(i) An irrelevant term has \(v_t=0\): its orbital integral contains
the unipotent average, which vanishes when the characters disagree, as
stated in the \hyperlink{proof.gram-orbit-expansion}{orbit expansion}.
The edition takes this vanishing of the irrelevant degenerate vectors
from the release
\releasecite{Formalism/LocalToGlobal/SiegelLatticeTheory/CokernelGLCounts/RankPointCounts/CovolumeSpans.lean}{35}{vanishing of the irrelevant degenerate vectors}.

(ii) A relevant term has \(1\le j(t)\le r_S\). Here \(j(t)\) is the
index \(j\) of the normal form, so \(\dim(X\cap X')=r_S-j(t)\) and the
upper bound holds by definition. Suppose \(j(t)=0\). Then
\(\dim(X\cap X')=r_S\), so \(X=X'\). By Step~2 every nonzero
\(x\in X\) has the same level for both flags, so \(X'_b=X_b\) for every
\(b\). Since \(r_S\ge1\), Step~2 also gives \(z'-\theta z\in X\) with
\(\theta\in E^\times\), so \(X'^+=X^+\). Since \(w=v\), \(z\perp X\) and
\(X\) is isotropic, comparing \(\langle z',z'\rangle\) with
\(\langle z,z\rangle\) gives \(\theta\bar\theta=1\). The release shows
that a relevant term whose flags have equal subspaces \(X_b\) and equal
\(X^+\) is the identity term, using that the word prime \(T_w\) is
prime to \(g_0g_1\)
\releasecite{Formalism/LocalToGlobal/SiegelLatticeTheory/CokernelGLCounts/RankPointCounts/IsIdentityOfEqualFlags.lean}{20}{identity recognition for equal flags}.
The argument there runs as follows. The comparison scalar is a unit
with all absolute values one, hence a root of unity. Admissibility of
the lattice data makes it congruent to \(1\) modulo the primes above
\(g_0\) and \(g_1\), hence equal to \(1\), and the term is then
realised by an element of the unipotent group of the lattice. The identity term is excluded from the
degenerate family, so \(j(t)\ge1\).

(iii) The typed real amplitude bound. Set \(H_c=b^{cm^2f}\) and
\(H_1=b^{C_Hm}\). A relevant term \(t\) has a label
\(\delta=(\delta_l)_l\), the local comparison data of its two flags at
the levels \(l\le j(t)\), with \(\delta_l=0\) for \(l>j(t)\). Let
\(\Delta(t)\) be the absolute norm of the local discriminant attached
to \(\delta_{j(t)}\). Put
\(S(t)=\prod_\tau\max\bigl(1,\max_l|\delta_l|_\tau/(2H_1)^{2r_S}\bigr)\ge1\),
where the product runs over the real places \(\tau\) of \(F\) and
\(|\delta_l|_\tau\) is the archimedean size of the label entry. The
release gives constants \(C_H,C_h,C_p\ge0\) with the following
property. For every \(C_s\ge C_p+4\) and every sufficiently large
\(C\ge0\), each finite set of same-row ordered pairs of relevant terms,
in which the first terms all have type \((j',\delta')\) and the second
terms all have type \((j'',\delta'')\), satisfies
\[
 \sum\|v_t\|\,\|v_{t'}\|\le\dim\pi_w\,M_w\,
   (ae)(j',\delta')\,(ae)(j'',\delta'').
\]
Here
\(a(j,\delta)=(C/k)^{j(k-1-j)f}\Delta^{-(k-1-j)/2}H_C^{Cr_S^4}S^{-C_sk}\)
and \(e(j,\delta)=\Delta^{k/2}H_{C_h}^{C_hr_S^4}S^{C_pk}\),
\(\pi_w\) is the fibre representation of the columns of \(w\), and
\(M_w=\sum_\Lambda|A_\Lambda|^{-1}\) is the total lattice weight
\releasecite{Formalism/LocalToGlobal/SiegelLatticeTheory/FlagStabilizerTensors/RowsPos.lean}{18}{typed real amplitude bound}.

(iv) \(0<\Delta(t)\le H_CS(t)^2\) for the same \(C\)
\releasecite{Formalism/LocalToGlobal/SiegelLatticeTheory/CokernelGLCounts/RankPointCounts/DimFormC.lean}{65}{discriminant bound}.

(v) There is \(C_c\ge0\) such that for each \(1\le j\le r_S\) and each
finite set \(T\) of labels realised by relevant terms of rank \(j\),
\(\sum_{\delta\in T}S^{-3r_S}\le H_{C_c}^{C_cr_S}\)
\releasecite{Formalism/LocalToGlobal/SiegelLatticeTheory/CokernelGLCounts/RankPointCounts/DimFormC.lean}{72}{coset count}.

(vi) The dimension formula \(\Gamma_1M_w\dim\pi_w\le d_w\)
\releasecite{Formalism/LocalToGlobal/SiegelLatticeTheory/CokernelGLCounts/RankPointCounts/DimFormC.lean}{80}{dimension formula}.

The constants are fixed in the order \(C_H,C_h,C_p\), then
\(C_s=C_p+4\), then \(C\) large for (iii) and (iv), then \(C_c\), and
finally \(C_A=1+C+C^2+C_h^2+C_c^2\). The assembly is now algebra. The
exponents of \(\Delta\) add to \((j+1)/2\), so
\[
 (ae)(j,\delta)=(C/k)^{j(k-1-j)f}H_C^{Cr_S^4}H_{C_h}^{C_hr_S^4}
   \Delta^{(j+1)/2}S^{(C_p-C_s)k}.
\]
By (iv), \(\Delta^{(j+1)/2}\le H_C^{(j+1)/2}S^{j+1}\). Since \(S\ge1\),
\(C_s-C_p\ge4\) and \(j+1+3r_S\le4r_S+1\le4k\), we get
\(S^{j+1+(C_p-C_s)k}\le S^{-3r_S}\). Summing over a finite set of
labels of rank \(j\) with (v) gives
\[
 \sum_\delta(ae)(j,\delta)\le\tilde\alpha_j=(C/k)^{j(k-1-j)f}
   H_C^{Cr_S^4+(j+1)/2}H_{C_h}^{C_hr_S^4}H_{C_c}^{C_cr_S}.
\]
Now take a finite set of same-row ordered pairs of nonidentity
degenerate terms. Pairs with an irrelevant member contribute \(0\) by
(i). By (ii), the rest split by \((j',\delta',j'',\delta'')\) with
\(1\le j',j''\le r_S\), and only finitely many labels occur. Applying
(iii) to each piece and summing over labels gives the bound
\(\dim\pi_w\,M_w(\sum_{j=1}^{r_S}\tilde\alpha_j)^2\). Absolute summability lets
the finite set exhaust the sum. Finally
\(H_c^y=\exp(cym^2f\log b)\), and \(r_S\le m\), \((j+1)/2\le m\) give
\(C^2r_S^4+C(j+1)/2+C_h^2r_S^4+C_c^2r_S\le C_Am^4\). Since
\(0\le C\le C_A\) and \(j(k-1-j)\ge0\), this gives \(\tilde\alpha_j\le\alpha_j\).
Multiplying by \(\Gamma_1\) and using (vi) gives the typed amplitude
bound.

The remaining numerics are as follows. Since \(m=2r_S+2\) and
\(k\ge m+3\), each \(1\le j\le r_S\) has \(j(k-1-j)\ge k-2\). Beyond a
scale threshold, \(\log k\ge8\log C_A\) and
\(k\log k\ge8C_Am^6\log b\). Indeed \(k\ge70\log b\) tends to
infinity, and \(k\ge8C_\Gamma m^5f\log b\ge8C_\Gamma m^9\log b\), so
the second inequality holds once \(C_\Gamma\log k\ge C_A\). Then
\(C_A/k\le1\), \(r_S\le k\), and
\[
 \log\Bigl(2\Bigl(\sum_j\alpha_j\Bigr)^2\Bigr)
 \le\log2+2\log k+2C_Am^6f\log b-2(k-2)f\log k+2(k-2)f\log C_A .
\]
The last two positive terms are each at most \(\tfrac14kf\log k\), and
\(\log2+2\log k\le3f\log k\). The right side is therefore at most
\((7-\tfrac32k)f\log k\le-kf\log k\), since \(k\ge14\). With
\(P_g\le2\) this gives \(\|K_{\rm deg}\|_{\rm HS}^2\le d_w\exp(-kf\log k)\).

\paragraph{The nondegenerate sum.}
\teamanchor{proof.equal-word-nondegenerate}%
The \hyperlink{proof.quadratic-row-contraction}{quadratic row contraction},
the \hyperlink{proof.orbital-error-estimates}{orbital error estimates}
and the combination in the
\hyperlink{proof.mixed-word-estimate}{mixed-word estimate} apply with
\(v=w\), with one change: the scalar main term. Put
\(D_{ww}=d_w^2/D_\star\). The edition takes the following equal-word
scalar bound from the release
\releasecite{Development/AbelianVarietyTheory/Proofs/FlagMeasures/RowMainEqPOfSiegelWeil.lean}{10}{Siegel--Weil}.
The signed sum over the generic scalar pairs of \(K_{ww}\), weighted by
\(\Gamma_1\) as in the finite convolution, is at most
\(\Gamma_BD_{ww}\), with \(\Gamma_B=\exp(C_Bm^5f\log b)\) for a
constant \(C_B\ge0\) fixed on the family. The release does not prove
this by the positivity and Poisson argument of the
\hyperlink{proof.scalar-row-normalization}{scalar normalization}, and
the edition does not use the bound \(AD_{wv}\) of that argument with
\(v=w\). Instead, the release groups the generic scalar pairs into
classes and, within a class, by a rational point and a normal datum.
It bounds each class by Cauchy--Schwarz and a mean bound for the local
generic amplitudes, and counts the points and normals with height
weights
\releasecite{Formalism/LocalToGlobal/SiegelLatticeTheory/CokernelGLCounts/RankPointCounts/MainEqCut.lean}{195}{class, amplitude and height bound for generic scalar pairs}.
Only then is the sum normalized. The Siegel--Weil formula, with the
local Bessel integral from Section~\ref{sec:periods} in both word slots
of the real orbital integral, enters through the equal-word mass data.
These mass data bound the \(\Gamma_1\)-weighted archimedean mass by
the scale \(d_wd_v/D_\star\) of the distinct-word normalization, now
taken at \(v=w\), times a factor that depends only on the field
discriminant, the auxiliary primes and the rank
\releasecite{Formalism/LocalToGlobal/SiegelLatticeTheory/CokernelGLCounts/RankPointCounts/MainEqCut.lean}{259}{normalization of the class bound by the mass data}.
The scale \(D_{ww}\) is the only thing the edition takes from the
scalar normalization at \(v=w\). This factor and the class and height
counts are absorbed into \(\Gamma_B\), which is why the bound is crude. It
suffices, because \(\Gamma\ge\Gamma_Bb^{m^5f}\) below. The error bounds
\(\epsilon_LD_{ww}\) and \(\epsilon_ND_{ww}\) are also taken for
\(v=w\); their release statements cover every pair of words.
The identity orbit has \(G=0\), so it lies in neither family. As in
the mixed-word estimate,
\[
 \|K_{\rm nd}\|_{\rm HS}^2
 \le D_{ww}\Bigl[\sqrt{P_g(\Gamma_B+2\epsilon_L)}
    +\sqrt{P_g(\epsilon_L+\epsilon_N)}\Bigr]^2 .
\]
The release
\releasecite{Formalism/LocalToGlobal/SiegelLatticeTheory/FlagStabilizerTensors/RowsPos.lean}{210}{combines these equal-word pair estimates}
in the same way. To bound the bracket put
\(E=\exp(\max(C_B,C_N)m^5f\log b)\ge1\). Then \(\Gamma_B\le E\),
\(\epsilon_L\le1\), and \(\epsilon_N\le E\) because
\(k\ge2r_S+5\). Using \((\sqrt a+\sqrt c)^2\le2(a+c)\) and \(P_g\le2\),
the bracket is at most \(4(\Gamma_B+3\epsilon_L+\epsilon_N)\le20E\).
The constants \(C_B\) and \(C_N\) are fixed first, and the family
exponent is chosen with
\[
 C_\Gamma\ge\max(C_B,C_N)+1 .
\]
Then \(\Gamma\ge E\,b^{m^5f}\), and \(b^{m^5f}\ge16^{512}>20\) because
\(b\ge16\), \(m\ge2\) and \(f\ge m^4\). Hence
\(\|K_{\rm nd}\|_{\rm HS}^2\le\Gamma d_w^2/D_\star\).
The rank thresholds are taken after \(C_\Gamma\), and the rank
satisfies \(k\ge 8C_\Gamma m^5f\log b\).

The upper ratio in \eqref{eq:trace-balance} gives
\(\Gamma d_w^2/D_\star\le(C_0/\Gamma)d_w\).
Combining the degenerate and nondegenerate estimates
by the squared triangle inequality
\(\|A+B\|^2\le2\|A\|^2+2\|B\|^2\) gives
\begin{equation}
 \label{eq:trace-single-diagonal}
 \|K_{ww}-I\|_{\mathrm{HS}}^2\le e_0 d_w,\qquad
 e_0=2\exp(-kf\log k)+2C_0/\Gamma .
\end{equation}
Here \(k\ge 8\log\Gamma\) and \(\log\Gamma\ge C_\Gamma x^\varepsilon\), so
\(\sqrt{e_0}\le\exp(-x^{\varepsilon/8})\) for large \(x\), uniformly
over the selected translates.

\noindent\leavevmode\teamanchor{candidate.diagonal}%
\textbf{AI-proposed candidate (the same-word Gram defect).}
Fix the target, precision and tied family with its levels, weight
gaps and rank conditions, including \eqref{eq:trace-balance} and
\(k\ge8\log\Gamma\), with \(C_\Gamma\ge\max(C_B,C_N)+1\) for the
fixed constants \(C_B\) of the equal-word scalar bound and \(C_N\) of
the nongeneric estimate. At every selected translate beyond one common
scale threshold, the full Gram operator of each word has mean
squared defect
\[
 \frac{\|L_w^*L_w-I\|_{\rm HS}^2}{d_w}
 \le 2e^{-kf\log k}+\frac{2C_0}{\Gamma},
 \qquad \Gamma=\exp(C_\Gamma m^5 f\log b).
\]
The normalization gives \(I\) on the identity double coset. The
estimate uses absolute bounds for the remaining degenerate terms
and the equal-word nondegenerate bound \(\Gamma d_w^2/D_\star\),
\hyperlink{proof.full-diagonal-estimate}{derived above} from the
release inputs named there.
This nondegenerate calculation shares orbital estimates with the
mixed-word proof below. The same-word statement controls the mean
of \((\sigma_j^2-1)^2\) over the \(d_w\) singular values of \(L_w\),
including zeros. Mixed bounds alone allow mutually orthogonal
images with each \(L_w\) twice an isometric inclusion, for which
the mean squared defect is \(9\).

Wei obtains a growing matrix of \(qI_2\)-Fourier coefficients of
the \(pI_2\) Siegel Poincar\'e series, with main term \(8I\) and
an entrywise error decreasing faster than any power of the weight.
\cite[Lemma~2.3 and \S3.1]{tr3wei-poincare}
The estimate above treats all full Bessel columns of one word
in the unitary family, with the indicated dependence on its rank,
field and level.

\leavevmode\teamanchor{proof.ramified-column-loss}%
The period source \(H_w^{\mathrm{act}}=\mathcal M_w\) is the ramified
column space included by \(\jmath_w\) in Section~\ref{sec:periods}.
At level zero the columns are sections over the definite class set,
with one residual compact representation fibre per class. Since \(T_w\)
splits completely in \(E\), there are \(f\) places of \(F\) over it
and the \(T_w\)-diamond group is
\(\mathsf A_w=\prod_{z\mid T_w}\mathbf F_{T_w}^{\times}\).
It acts by right translation on the classes and carries the fibres
isomorphically along this action. A diamond representative is an
integral local dilation at the \(T_w\)-places and the identity at
the auxiliary places. If it fixes a class, the double-coset equation
identifies its product with a level element with a rational point
preserving the lattice of that class. This rational stabilizer is
finite by definiteness. At each auxiliary prime the reduction is
upper unitriangular, and the preimage of that group is pro-\(g_i\)
because the principal congruence kernel is pro-\(g_i\). Thus the
stabilizer element has order a power of each of two distinct primes.
It is the identity, so the diamond representative is itself in the
level and acts by the trivial multiplier at every \(T_w\)-place.
For comparison, the release proves the same freeness at
\releasecite{Formalism/LocalToGlobal/GlobalTheory/GL6/DiagonalSlopes/EqOneOfCoprimeTorsionTwoPrimes.lean}{670}{Freeness at the level-zero classes}.

On each free orbit fix an identification of the fibre at one class
and transport it by the group action. This identifies the full
column representation over the orbit with the regular representation
of \(\mathsf A_w\) tensored with that fibre. In each
\(\mathbf C[\mathbf F_{T_w}^{\times}]\) the kernel of the local
averaging projector has dimension \(T_w-2\). The actual ramified
space is the intersection of these local kernels. Its restriction
to the orbit is the tensor product of the individual kernels with
the fibre, since the \(f\) local projections act in distinct factors.
Adding the ranks over the orbits therefore gives the first equality
below. With \(r_w=\dim\mathcal M_w\) and \(T_w\ge b\ge3\), Bernoulli's
inequality gives the rest:
\begin{equation}
 \label{eq:trace-ramification}
 r_w=d_w\left(1-\frac1{T_w-1}\right)^f,\qquad
 0\le d_w-r_w\le\frac{2f}{b}d_w.
\end{equation}
This is the ratio for the same level-zero columns counted in
\eqref{eq:trace-balance}.
For comparison, the release states the same formula at
\releasecite{Formalism/SharpRestriction/GlobalSchurBounds/FiniteModelOfAdmissibleTransports.lean}{842}{Ramified rank formula}.

The diagonal input gives another use of this calculation.
On \(\ker\mathcal P_w\), the operator \(K_{ww}-I\) is \(-I\).
Consequently \(\dim\ker\mathcal P_w\le e_0 d_w\). The bound
\(f\le k\le x^5\) implies \(e_0<1-2f/b\) for sufficiently large \(x\),
so some vector in \(\mathcal M_w\) has nonzero image under
\(\mathcal P_w\). Pairing with that image in the full \(\mathcal H\)
supplies the nonzero bottom pairing.

In the construction of the bottom reference pairing, the
\releasecite{Development/AbelianVarietyTheory/Proofs/DecoratedLatticeProofs/DiagonalRowOfSiegelWeilUnitary.lean}{50}{earlier single-word estimate}
comes from the same local Bessel integrals and Siegel--Weil. In
combination with the ramified rank bound, it puts the squared diagonal
defect strictly below the dimension of the ramified column space of
each such bundle. The
\releasecite{Development/AutomorphicArithmetic/ThetaCorrespondence/ThetaPairLifts/ThirdPlaceSeries/MainTheoremPublishedResultsFinal.lean}{34}{formal construction}
uses this inequality to obtain a nonzero image of a ramified column.
For the conditional argument we use the period maps supplied by
\hyperlink{candidate.periods}{the stated comparison}.

\subsection{Squaring the orbital sum at finite level}
\leavevmode\teamanchor{proof.quadratic-row-contraction}%
Fix \(w\ne v\) and let \(\mathcal S\) be either of the two orbit
families with \(G\ne0\). The same grouping applies with \(w=v\). The
identity orbit has \(G=0\) and lies in neither family. The resulting
equal-word bound \(\Gamma d_w^2/D_\star\) is derived in the
\hyperlink{proof.equal-word-nondegenerate}{single-word estimate} from
the release inputs named there. Group the terms by their two lattice
representatives and the remaining index \(i(t)\) in the finite Levi
group. Call this triple the row \(r(t)\). At a place of \(F\) above
\(g_0g_1\) with residue size \(q\), the group is
\(G_q=\mathrm{GL}_k(\mathbf F_q)\) with upper unitriangular subgroup
\(U_q\). Write \(I=\prod_q G_q\), with all such products indexed by
these places. The weighted vector \(v_t\) is the normalized real
orbital operator in the Hilbert--Schmidt summand for its lattice pair,
multiplied by the finite unipotent character and the square root of
the product of the two lattice weights.
Put \(K_{\mathcal S}=\sum_{t\in\mathcal S}A_t\) and
\(V_i=\sum_{t\in\mathcal S:\,i(t)=i}v_t\).
In the chosen column representatives the unipotent factor of the
finite character is constant for each lattice pair. The remaining
factor is the character of \(\prod_q U_q\), extended by zero, evaluated
at the relative Levi index for the two column representatives. The
representatives give free decompositions of \(I\), and the Gram sums
for each family are invariant under the corresponding left and right
translations in \(I\). The edition takes these structural facts
from the release:
\releasecite{Formalism/LocalToGlobal/SiegelLatticeTheory/ClosedLocalOrbits/GaugeLatticeShifts/ColPhaseCut.lean}{49}{the character factorization},
\releasecite{Formalism/LocalToGlobal/SiegelLatticeTheory/ClosedLocalOrbits/GaugeLatticeShifts/ChiNCut.lean}{48}{the constant unipotent factor},
\releasecite{Formalism/LocalToGlobal/SiegelLatticeTheory/ClosedLocalOrbits/GaugeLatticeShifts/NormExpSatHolds.lean}{17}{the free column decompositions}
and the left and right Gram invariance
\releasecite{Formalism/LocalToGlobal/SiegelLatticeTheory/ClosedLocalOrbits/GaugeLatticeShifts/NormExpSatHolds.lean}{23}{left invariance}
\releasecite{Formalism/LocalToGlobal/SiegelLatticeTheory/ClosedLocalOrbits/GaugeLatticeShifts/NormExpSatHolds.lean}{37}{right invariance}.
The derivation below shows that they give the squared norm
\[
  \|K_{\mathcal S}\|_{\mathrm{HS}}^2
    =\operatorname{Re}\sum_{i,j\in I}\mathcal G_{ij}\langle V_i,V_j\rangle,
  \qquad
  \mathcal G_{ij}=\overline{\prod_q\gamma_q((j^{-1}i)_q)}.
\]
Here \(\gamma_q=\operatorname{ch}(\operatorname{Ind}_{U_q}^{G_q}\psi_q)\)
is the Gelfand--Graev character,
where \(\psi_q\) is the nondegenerate character obtained from the
chosen nontrivial residue-field character on the sum of the simple
superdiagonal entries.

\leavevmode\teamanchor{proof.squared-norm-identity}%
\paragraph{Derivation of the squared norm.}
Take the Hilbert--Schmidt inner product conjugate-linear in the first
variable. Let \(U=\prod_qU_q\), let \(\psi=\prod_q\psi_q\) on \(U\),
extended by zero to \(I\), and let \(\gamma=\prod_q\gamma_q\). Then
\(\psi(ug)=\psi(u)\psi(g)\) and \(\psi(gu)=\psi(g)\psi(u)\) for \(u\in U\)
and \(g\in I\). Fix a
lattice pair \(\Lambda\) occurring in \(\mathcal S\). The structural
facts say the following. The \(\Lambda\)-summand of \(K_{\mathcal S}\)
has, between a column \(\rho\) of the first lattice and a column
\(\rho'\) of the second, the block
\(\sum_t\psi(s_\rho^{-1}i(t)s'_{\rho'})a_t\). The sum runs over
\(t\in\mathcal S\) with lattice pair \(\Lambda\), \(s_\rho,s'_{\rho'}\in
I\) are the column representatives, and \(a_t\) is the remaining factor of \(A_t\): the normalized real
orbital operator times the level volume and the unipotent factor,
which is constant for the lattice pair.
There are subsets \(A_\Lambda,A'_\Lambda\subseteq I\) for which
\((b,\rho,u)\mapsto bs_\rho u\) and \((b,\rho',u)\mapsto bs'_{\rho'}u\)
are bijections onto \(I\) from \(A_\Lambda\times\{\rho\}\times U\) and
\(A'_\Lambda\times\{\rho'\}\times U\), where \(\rho\) and \(\rho'\) run
over the columns of the first and second lattice of \(\Lambda\). The two lattice
weights are \(|A_\Lambda|^{-1}\) and \(|A'_\Lambda|^{-1}\). Finally,
\(Z_\Lambda(i)=\sum_{t:\,i(t)=i}a_t\) satisfies
\(\langle Z_\Lambda(bi),Z_\Lambda(bj)\rangle
=\langle Z_\Lambda(i),Z_\Lambda(j)\rangle\) for \(b\in A_\Lambda\) and
\(\langle Z_\Lambda(ib^{-1}),Z_\Lambda(jb^{-1})\rangle
=\langle Z_\Lambda(i),Z_\Lambda(j)\rangle\) for \(b\in A'_\Lambda\).
The sums over \(t\) converge absolutely in the Hilbert--Schmidt norm
\releasecite{Formalism/LocalToGlobal/SiegelLatticeTheory/ClosedLocalOrbits/GaugeLatticeShifts/KExpCut.lean}{301}{absolute summability of the orbit terms},
which justifies the regrouping below.
In the release this summability is deduced from three facts: each
block of a term has norm at most a fixed volume factor times the
operator norm of the orbital integral of that term
\releasecite{Formalism/LocalToGlobal/SiegelLatticeTheory/ClosedLocalOrbits/GaugeLatticeShifts/KExpCut.lean}{95}{block bound},
this operator norm is at most the integral, over the fundamental
domain of the term at the real places, of the product of the norms of
the two coefficient vectors \(W_w\) and \(W_v\)
\releasecite{Formalism/LocalToGlobal/SiegelLatticeTheory/ClosedLocalOrbits/GaugeLatticeShifts/DoubleCosetSingleUnfolding.lean}{383}{norm of an orbital integral},
and these absolute integrals have a finite sum over the terms
\releasecite{Formalism/LocalToGlobal/SiegelLatticeTheory/ClosedVolumeStatements/FullExistsTiedHolds.lean}{27}{absolute convergence over the terms}.
The terms there are indexed by lattice pairs and adelic classes modulo
the lattice unipotent subgroups. This indexing refines the rational
double cosets of the unfolding, so the edition does not deduce the last
finiteness from the absolute convergence proved for the unfolding.
The Hilbert--Schmidt norm is additive over blocks, so
\[
  \|K_{\mathcal S}\|_{\mathrm{HS}}^2=\sum_\Lambda\sum_{\rho,\rho'}
   \Phi_\Lambda(s_\rho,s'_{\rho'}),\qquad
  \Phi_\Lambda(x,y)=\sum_{i,j\in I}
   \overline{\psi(x^{-1}iy)}\,\psi(x^{-1}jy)\,
   \langle Z_\Lambda(i),Z_\Lambda(j)\rangle .
\]
Replacing \(x\) by \(xu\) or \(y\) by \(yu\) with \(u\in U\) multiplies
the summand by \(|\psi(u)|^2=1\). Replacing \(x\) by \(bx\) with
\(b\in A_\Lambda\), or \(y\) by \(by\) with \(b\in A'_\Lambda\), is
undone by substituting \(i\mapsto bi\), \(j\mapsto bj\), or
\(i\mapsto ib^{-1}\), \(j\mapsto jb^{-1}\), and using the Gram
invariance. By the free decompositions,
\[
  \sum_{\rho,\rho'}\Phi_\Lambda(s_\rho,s'_{\rho'})
   =\frac{1}{|A_\Lambda||A'_\Lambda||U|^2}\sum_{x,y\in I}
     \Phi_\Lambda(x,y).
\]
In the sum over \(y\), substitute \(y\mapsto i^{-1}xy\). The term
\(\overline{\psi(y)}\psi(x^{-1}ji^{-1}xy)\) vanishes unless \(y\in U\)
and \(x^{-1}ji^{-1}x\in U\), and then it equals
\(\psi(x^{-1}ji^{-1}x)\). The Frobenius formula
\(\gamma(g)=|U|^{-1}\sum_{x\in I}\psi(x^{-1}gx)\) then gives
\(\sum_{x,y}\overline{\psi(x^{-1}iy)}\psi(x^{-1}jy)=|U|^2\gamma(ji^{-1})\).
Since \(\gamma\) is a class function and \(\gamma(g^{-1})
=\overline{\gamma(g)}\), we have \(\gamma(ji^{-1})=\gamma(i^{-1}j)
=\mathcal G_{ij}\). Hence
\[
  \|K_{\mathcal S}\|_{\mathrm{HS}}^2=\sum_{i,j}\mathcal G_{ij}
   \sum_\Lambda\frac{\langle Z_\Lambda(i),Z_\Lambda(j)\rangle}
     {|A_\Lambda||A'_\Lambda|}
   =\sum_{i,j}\mathcal G_{ij}\langle V_i,V_j\rangle ,
\]
because \(V_i\) has \(\Lambda\)-component
\((|A_\Lambda||A'_\Lambda|)^{-1/2}Z_\Lambda(i)\) in the orthogonal sum
over lattice pairs of the Hilbert--Schmidt summands. The left side is
real, so taking real parts changes nothing. Different lattice pairs
are orthogonal by this construction.

Let \(\Gamma_1=\prod_q\gamma_q(1)=\prod_q[G_q:U_q]\)
and \(P_g=\prod_q(1+3/q)\).
The finite convolution gives
\[
  0\le\mathcal G\le P_g\Gamma_1\operatorname{Id}_{\mathbf C^I},\qquad
  \|K_{\mathcal S}\|_{\mathrm{HS}}^2
  \le P_g\Gamma_1\sum_i\|V_i\|^2
  =P_g\Gamma_1\!
     \sum_{\substack{t,t'\in\mathcal S\\r(t)=r(t')}}
          \operatorname{Re}\langle v_t,v_{t'}\rangle.
\]
For the inequality, use an orthonormal basis of the span of the
finitely many \(V_i\) and apply the scalar matrix bound to each
coordinate vector. For the last equality, within each \(i\) different
lattice pairs have orthogonal summands. Absolute convergence then
permits the expansion over ordered pairs in each row.

The character \(\overline{\psi_q}\) is conjugate to \(\psi_q\) by a
diagonal matrix whose successive diagonal entries have ratio \(-1\).
Consequently \(\gamma_q\) is real. For convolution with counting
measure, character orthogonality and
Frobenius reciprocity give the eigenvalue of
\(\overline{\gamma_q}/\gamma_q(1)\) on the irreducible \(\rho\) as
\[
  \frac{|U_q|}{\dim\rho}\,
        \dim\operatorname{Hom}_{U_q}(\psi_q,\rho).
\]
For \(\mathrm{GL}_k(\mathbf F_q)\) over a finite field and the
nondegenerate \(\psi_q\) just defined, Gel'fand proves that the complex
representation \(\operatorname{Ind}_{U_q}^{G_q}\psi_q\) is
multiplicity-free. For the mirabolic \(P_q\), defined by
\(e_k^{\mathsf T}g=e_k^{\mathsf T}\) with \(e_k\) the last basis vector,
\(\operatorname{Ind}_{U_q}^{P_q}\psi_q\) is irreducible of
dimension \([P_q:U_q]=\prod_{j=1}^{k-1}(q^j-1)\).
Thus for each irreducible complex \(\rho\) of \(G_q\), the displayed
\(\psi_q\)-multiplicity is \(0\) or \(1\). Frobenius reciprocity for
\(U_q\subset P_q\) gives exactly one copy of the mirabolic induction
in \(\rho|_{P_q}\) when the multiplicity is \(1\), and therefore gives
the same lower bound for \(\dim\rho\).
\cite[Theorem~0.5 and Proposition~1.2, pp.~17--18]{tr3gelfand}
With \(|U_q|=q^{k(k-1)/2}\), the local eigenvalue
is at most \(\prod_{j=1}^{k-1}(1-q^{-j})^{-1}\le 1+3/q\) for \(q\ge3\).
The product before inversion is at least \(1-1/(q-1)\).
Taking tensor products over the places yields the bound for \(\mathcal G\).

For a pair of terms in the same row, consider the fixed flag
\(\mathfrak f\) and two moving
flags \(\mathfrak f',\mathfrak f''\), set \(W_1\) and \(W_2\) equal to
the orthogonal complements of the two joint augmented frames. Let
\(\beta\) be the isometry comparing the moving flags and let \(X\)
and \(X''\) be the fixed and second moving isotropic spans. The comparison
space is the subspace of \(x\in W_1\) for which there exists
\(y\in W_2\) with \(x-y\in X\) and \(\beta x-y\in X''\).
The corank \(s_0\) is \(k\) minus its dimension, so \(s_0\ge1\)
because \(\dim W_1=k-1\). To define a scalar pair, write
\(U' = X^+(\mathfrak f)+X^+(\mathfrak f')\) for the joint span of the
first pair of augmented frames, and \(U''\) for the analogous span for
the second. In \(U'\), take lifts \(y\) and \(y'\) perpendicular to
\(X^+(\mathfrak f)\) and \(X^+(\mathfrak f')\), respectively. If
\(y-y'=\lambda z+\lambda' z'\) modulo the sum of the two isotropic
spans, call \((\lambda,\lambda')\) a pair of distinguished-vector
coefficients. Compare this
relation with the corresponding one for lifts \(\tilde y,\tilde y'\)
in \(U''\) such that \(y-\tilde y\in X\) and
\(\tilde y'-\beta y'\in X''\). A pair is scalar if it has corank one
and the
\releasecite{Formalism/SharpRestriction/CuspSpanNormBounds/IsPrimitiveSub.lean}{252}{two coefficient relations agree}
for every such choice of lifts.

\subsection{The leading contribution to the mixed Gram blocks}

\leavevmode\teamanchor{proof.trace-word-ratios}%
Keep \(w\ne v\) fixed and set \(D_{wv}=d_wd_v/D_\star\).
The word primes \(T_w\) split completely in \(E\) and are congruent
to \(1\) modulo \((g_0g_1)^2\). We compare them, and the compact
Vandermondes, for two words. Write \(j_0=k-2\).
At a real place, the compact word parameters have two common
endpoints and the \(j_0\) interior entries \(v_i-1/2-\Delta b_i\),
where \(b_i\in\{0,1\}\) and \(v_i-v_{i+1}=4\Delta\).
The upper endpoint \(A\) satisfies \(A-v_1\ge4k\Delta\).
The lower endpoint is \(v_{k-1}-1/2\) away from the designated top
places, and \(-b_{\rm top}\) with \(b_{\rm top}\ge B_1/2\) at
the top places. The selected weight data give
\(\Delta\ge1\), \(v_i\ge0\) and \(B_1\ge4kx\Delta\).
\releasecite{Formalism/SharpRestriction/CuspSpanNormBounds/BruhatCellsKronecker/CfgTopFalseIff.lean}{39}{Endpoints and the shifted grid}
\releasecite{Formalism/SharpRestriction/CuspSpanNormBounds/StandingNumericalData.lean}{84}{Endpoint inequalities}
\releasecite{Formalism/SharpRestriction/GlobalSchurBounds/StandardLattice.lean}{500}{Weight size}.

For a single word the interior log Vandermonde relative to the
unshifted grid is
\[
 \sum_{h=1}^{j_0-1}\sum_{i=1}^{j_0-h}
    \log\left(1+\frac{b_{i+h}-b_i}{4h}\right).
\]
The linear term for lag \(h\) is bounded in absolute value by
\(\min(h,j_0-h)/(4h)\), by cancellation of the overlapping sums of
bits. For \(|t|\le1/4\), integration of \(1/(1+t)\) gives
\(|\log(1+t)-t|\le(2/3)t^2\). The sum of these remainders is at
most \((j_0/24)\sum_{h\ge1}h^{-2}\le j_0/12\) by integral
comparison. The difference for two words on the interior is
therefore at most \(2j_0/3\) in absolute value. At the upper endpoint the
log change of each gap is bounded by \(\Delta/(4k\Delta)\).
At the lower endpoint outside the top places the gap for interior
entry \(i\) is at least \(3(k-1-i)\Delta\), so the total log change
there is at most \(\sum_{j=1}^{j_0}1/(3j)\). At a top place the
interior entries are at least \(-\Delta-1/2\), so each gap to
the lower endpoint is at least \(k\Delta\) for large \(x\), at
total cost at most \(1\). The difference of the two endpoints is
fixed. For large \(k\) the resulting logarithmic difference of the
full compact Vandermondes is at most \(3k/4\), at every real place.

The compact Weyl dimension formula gives \(d_w^{(0)}\) as a common
finite-level class factor times the product over real places of
the compact Vandermondes, divided by the same Weyl denominators.
This is \(d_w^*\) in the compact dimension formula of
Section~\ref{sec:integral}. Hence
\(\big|\log(d_w^{(0)}/d_v^{(0)})\big|\le3kf/4\).
Thus \(|\log(Y_w/Y_v)|\le3/4\). By the chosen prime interval,
every \(\log(T_w/Y_w)\) lies in \([0,\log(1+1/(fk))]\).
Combining the bounds gives
\begin{equation}
\label{eq:trace-word-ratios}
  |\log(T_w/T_v)|\le \tfrac34+\log(1+1/(fk))\le1,\qquad
  \left|\log\frac{\operatorname{Vand}_{z}(w)}
                   {\operatorname{Vand}_{z}(v)}\right|\le \tfrac34 k
        \quad(z\mid\infty).
\end{equation}
Here \(\operatorname{Vand}_z(w)\) is the compact Vandermonde at \(z\).
Thus \(T_w/T_v\) and the \(k\)-th roots of the compact Vandermonde
ratios lie in \([e^{-1},e]\).
These two-word comparisons are not inputs of the mixed estimate. Among
the release's size hypotheses for it
\releasecite{Formalism/LocalToGlobal/SiegelLatticeTheory/FlagStabilizerTensors/TheoremSRecursion.lean}{113}{size hypotheses of the mixed estimate},
\eqref{eq:trace-balance}, the column counts and the lower bounds on
\(T_w\) and \(T_v\) are imposed for each word separately, and no
comparison between the two words is assumed. The comparison of
critical levels is used in the mass-ratio step of the nonvanishing
argument
\releasecite{Formalism/SharpRestriction/CuspSpanNormBounds/BruhatCellsKronecker/FamilyCriticalLevelLowerBound.lean}{34}{critical-level ratio}.
The orbital estimates use \eqref{eq:trace-balance} together with
the prescribed gaps between archimedean weights and the local lattice
models. These hypotheses allow ai agent to
\releasecite{Formalism/LocalToGlobal/SiegelLatticeTheory/FlagStabilizerTensors/LinkSFinalP.lean}{105}{apply the lattice-counting and orbital-decay estimates to each pair of words}.

\leavevmode\teamanchor{proof.primitive-frame-densities}%
The leading term comes from the generic scalar pairs of corank one.
At an ordinary finite place \(\nu\), write \(E_\nu=E\otimes_F F_\nu\),
assumed split or unramified quadratic over \(F_\nu\). Let \(\varpi\)
be a uniformizer of \(F_\nu\), \(q=|\mathcal O_{F_\nu}/\varpi|\),
and \(J_\nu\) the Gram matrix in an integral basis of a self-dual
Hermitian lattice of rank \(r\). For a nonsingular integral Hermitian
matrix \(S\) of size \(\ell\), let \(c\) be the corank of its reduction.
At congruence depth \(n\), count matrices
\(X\in M_{r,\ell}(\mathcal O_{E_\nu}/\varpi^n)\) with
\(X^*J_\nu X\equiv S\pmod{\varpi^n}\) and with injective reduction
modulo \(\varpi\). In the split case injectivity is required in
each of the two factors. For \(r>\ell\) and \(\ell+c\le r\), the
limit of this count divided by \(q^{n(2r\ell-\ell^2)}\) is
\[
  \prod_{j=r-\ell-c+1}^{r}(1-\xi^j q^{-j}),\qquad
  \xi=\begin{cases}1 &\text{at a split place},\\ -1 &\text{at an inert place}.
       \end{cases}
\]
For the unramified quadratic field \(E_\nu\), the self-dual lattice
has Gram matrix \(I_r\) in an integral orthonormal basis. This is
Hironaka's primitive density for that ambient Gram matrix. The \(c\)
positive elementary-divisor exponents of \(S\) give denominator
\(\prod_{j=1}^{r-\ell-c}(1-(-q)^{-j})\) in her quotient of products.
\cite[Equation~(3.1), p.~572, and Corollary~3.3(i), p.~576]{tr3hironaka}
At a split place, let \(X_1,X_2\) denote the two reduced matrices after
identifying the ambient form with the standard one. In
\(X_1^{\mathsf T}X_2=\bar S\), with \(\bar S\) the reduction in the
first component, choosing \(X_1\) gives factors
\(j=r-\ell+1,\ldots,r\). The free block of \(X_2\) must be injective on
\(\ker\bar S\), giving the \(c\) remaining factors. Surjectivity of
\(\dot X_2\mapsto X_1^{\mathsf T}\dot X_2\) gives the same normalized count
at every congruence level.

The same product holds for a singular integral Hermitian \(S\), with
\(c\) the corank of \(\bar S\). If \(X\) has injective reduction, then
\(X^*J_\nu\) is surjective modulo \(\varpi\). Every Hermitian \(H\) is
\(T+T^*\) for some \(T\in M_\ell(\mathcal O_{E_\nu})\), because the trace
of the unramified or split algebra \(E_\nu\) is onto \(\mathcal O_{F_\nu}\).
Hence the differential \(\dot X\mapsto X^*J_\nu\dot X+\dot X^*J_\nu X\)
is onto the Hermitian matrices modulo \(\varpi\). Each solution modulo
\(\varpi^n\) with injective reduction therefore has exactly
\(q^{2r\ell-\ell^2}\) lifts modulo \(\varpi^{n+1}\), and the normalized
count equals its value at \(n=1\). That value depends only on
\(\bar S\). A nonsingular \(S'\equiv S\pmod\varpi\) has the same
reduction, so the displayed product applies. Such an \(S'\) is
\(S+\varpi tI_\ell\) for a suitable \(t\in\mathcal O_{F_\nu}\), since its
determinant is a nonzero polynomial in \(t\). This covers the
single-frame Gram matrix \(\operatorname{diag}(0_{r_S},-T_w)\), with
\(c=r_S\) at places prime to \(T_w\).

\leavevmode\teamanchor{proof.t3b-depth-one-count}%
The value at \(n=1\) can also be computed directly. The computation
uses no hypothesis on the residue characteristic, so it covers
\(q=2\) as well, and it gives the inert product independently of
Hironaka's formula. By \eqref{eq:classical-host-inertia}, \(E\) is unramified over \(F\)
at every finite place. The computation therefore applies at every
finite place of \(F\) at which the ambient lattice is self-dual,
including those dividing \(T_wT_v\) and those above the auxiliary
primes. Self-duality at every finite place is part of the
identification taken from the release below.
Let \(k_E\cong\mathbf F_{q^2}\) be the residue field of \(E_\nu\) at an
inert place, and \(h\) the reduction of \(J_\nu\), a nondegenerate
Hermitian form on \(k_E^r\). A Hermitian form \(h'\) with
\(h'(x,x)=0\) for all \(x\) is zero: expanding \(h'(x+ay,x+ay)=0\) gives
\(\operatorname{Tr}_{k_E/\mathbf F_q}(\bar a\,h'(x,y))=0\) for all
\(a\in k_E\), and the trace is a nonzero linear form. Hence every
Hermitian form over \(k_E\) is diagonalizable in every characteristic,
and \(\bar S=B^*\operatorname{diag}(a_1,\ldots,a_n,0_c)B\) with \(B\)
invertible, \(n=\ell-c\) and \(a_i\in\mathbf F_q^\times\). The map
\(X\mapsto XB^{-1}\) is a bijection between the injective solutions for
\(\bar S\) and those for the diagonal matrix. In a nondegenerate space
of dimension \(N\), let \(Z(N)\) be the number of vectors of norm \(0\),
including \(0\), and \(M(N)\) the number of vectors of a given norm
\(a\ne0\). The norm \(k_E^\times\to\mathbf F_q^\times\) is onto with
fibres of size \(q+1\), so \(M(N)\) does not depend on \(a\). Splitting
off an orthogonal line of norm one gives \(Z(1)=1\), \(M(1)=q+1\) and
\[
 Z(N)=Z(N-1)+(q^2-1)M(N-1),\qquad
 M(N)=(q+1)Z(N-1)+(q^2-q-1)M(N-1).
\]
By induction,
\(M(N)=q^{N-1}(q^N-(-1)^N)\) and
\(Z(N)-1=(q^N-(-1)^N)(q^{N-1}-(-1)^{N-1})\).
Now choose the columns in order. For \(i\le n\) the column \(x_i\) has
norm \(a_i\) and is orthogonal to the nondegenerate span of
\(x_1,\ldots,x_{i-1}\). Its orthogonal complement is nondegenerate of
dimension \(r-i+1\), which gives \(M(r-i+1)\) choices. The remaining
columns lie in the orthogonal complement \(W\) of
\(x_1,\ldots,x_n\), which is nondegenerate of dimension \(N=r-n\).
They must be isotropic, mutually orthogonal and linearly independent.
Let \(Y\) be the span of the first \(i-1\) of them. Then
\(Y^\perp\cap W=Y\oplus U\) with \(U\) nondegenerate of dimension
\(N-2i+2\), and this is at least \(2\) because \(\ell+c\le r\). A vector
\(y+u\) with \(y\in Y\) and \(u\in U\) is admissible exactly when \(u\) is
a nonzero isotropic vector. This gives
\(q^{2(i-1)}(Z(N-2i+2)-1)\) choices for \(i=1,\ldots,c\). Write each
count as a power of \(q\) times factors \(1-(-q)^{-j}\). The first
\(n\) columns give the factors \(j=r-n+1,\ldots,r\). The isotropic
columns give \(j=r-\ell-c+1,\ldots,r-n\). The powers of \(q\) add up to
\(q^{2r\ell-\ell^2}\). Hence the normalized count at \(n=1\) is the
displayed product with \(\xi=-1\). The argument at split places given
above also uses no hypothesis on \(q\). The Hensel step uses only the
surjectivity of the trace, which holds in residue characteristic
\(2\) as well.

In the release,
\releasecite{Formalism/LocalToGlobal/SiegelLatticeTheory/CokernelGLCounts/RankPointCounts/DualFamilyCounting/DensAtD.lean}{58}{self-duality and the local Gram matrices}
identify the orbital measures with this normalization.

In the scalar calculation \(\ell=m=2(r_S+1)\) and
\(k\ge \ell+f+5\). The first exponent in the density is
\(k-c\ge k-\ell\). At ordinary places the reference density is the
product of the two single-frame factors. At the places above the
auxiliary primes it uses one common single-frame factor. The ratio
of the pair density to the reference factor is at most
\(\exp(4q^{-(k-\ell)})\) at each place. At a finite place where the local
Gram matrix of the pair is not integral there are no integral
solutions, and the local pair factor is \(0\). Otherwise, by the
identification just cited, it is the displayed product with \(\ell=m\) and \(c\le\ell\) the corank of the
reduction of that matrix. The condition \(\ell+c\le r\) holds because
\(2\ell\le r=k+\ell-1\). The first exponent is \(a=k-c\ge k-\ell\ge2\).
Put \(x=q^{-1}\le1/2\). For \(j\ge a\),
\(1-x^j\le1-\xi^jx^j\le1+x^j\le e^{x^j}\). Also \(1-t\ge e^{-2t}\) for
\(0\le t\le1/2\), and \(\sum_{j\ge a}x^j\le2x^a\). Every product
\(\prod_{j=a}^{r}(1-\xi^jx^j)\) therefore lies between \(e^{-4x^a}\) and
\(e^{2x^a}\). In particular the pair factor is at most
\(e^{2x^{k-\ell}}\). The reference factor is a product of one or two
single-frame factors for \(\operatorname{diag}(0_{r_S},-T)\), which has
size \(r_S+1\) and residual corank \(c_1\in\{r_S,r_S+1\}\). Their first
exponent is \(r-(r_S+1)-c_1+1=k+r_S+1-c_1\ge k\), so each factor is at
least \(e^{-4x^k}\). Since \(\ell=m\ge4\), \(x^k\le x^{k-\ell}/16\), and
the reference factor is at least \(e^{-x^{k-\ell}/2}\). The ratio is
therefore at most
\(e^{5x^{k-\ell}/2}\le\exp(4q^{-(k-\ell)})\).
Let \(s_1=k-\ell\ge2\). There are at most \(f\) places of \(F\) above
each rational prime, so
\(\sum_{\mathfrak p}(\mathrm N\mathfrak p)^{-s_1}\le
f\sum_{n\ge2}n^{-s_1}\le3f\,2^{-s_1}\).
The upper bound \(E_{\rm den}\) for the normalized density products,
enlarged to at least \(1\), satisfies
\(1\le E_{\rm den}\le\exp(12f\,2^{-(k-\ell)})\).

\leavevmode\teamanchor{proof.scalar-row-normalization}%
Use the convergent unitary Siegel--Weil formula with \(Y_0\) of
Hermitian rank \(r>2\ell\) over the quadratic CM extension \(E/F\).
Let \(G_\ell\) be the unitary group of the split skew-Hermitian space of
rank \(2\ell\), with Siegel parabolic \(P_\ell\). Fix a nontrivial
additive character \(\vartheta\) of \(\mathbf A_F/F\) and a Hecke
character \(\chi\) of \(E\) restricting to \(\epsilon_{E/F}^{\,r}\) on
\(\mathbf A_F^\times\), where \(\epsilon_{E/F}\) is the character of the
quadratic extension. Write \(\omega\) for the corresponding Weil
representation on the Schwartz space of
\((Y_0\otimes_F\mathbf A_F)^\ell\).
If \(J\) is the Gram matrix of \(Y_0\), the form \(2J\) in Liu's
convention gives moment matrix \(X^*JX\). In particular the Siegel
unipotent element \(n(b)\) acts by the multiplier
\(\vartheta(\operatorname{tr}(bX^*JX))\) for Hermitian \(b\).

For any Schwartz \(\Phi\), set \(s_0=(r-\ell)/2\) and use the
Siegel--Weil section \(f_{\Phi,s}(g)=(\omega(g)\Phi)(0)\lambda(g)^{s-s_0}\),
where \(\lambda(g)=|\det a|_{\mathbf A_E}\) for the Siegel Levi
coordinate \(a\) in an Iwasawa decomposition of \(g\).
The Eisenstein series and the theta integral below both converge
absolutely at \(s_0\), and for every \(g\in G_\ell(\mathbf A_F)\)
\[
 E_\Phi(s_0,g)
 :=\sum_{\gamma\in P_\ell(F)\backslash G_\ell(F)}
       f_{\Phi,s_0}(\gamma g)
 =\int_{[G_0]}\sum_{X\in Y_0^\ell}(\omega(g)\Phi)(h^{-1}X)\,dh,
\]
where \(dh\) has total mass \(1\) on
\([G_0]=G_0(F)\backslash G_0(\mathbf A_F)\).
\cite[Theorem~2.1(1) and the Weil-representation formulas, pp.~858--859]{tr3liu}

For any nonsingular Hermitian \(\ell\)-by-\(\ell\) matrix \(S\)
over \(E\), take the \(S\)-Fourier coefficient \(E_{\Phi,S}(s_0,1)\),
integrating over the rational
quotient of the Siegel unipotent with multiplier
\(\vartheta(-\operatorname{tr}(bS))\) and quotient volume \(1\).
Taking this coefficient of the absolutely convergent formula above
gives equality with the corresponding coefficient of the theta
integral. Yamana gives this nonsingular coefficient identity for
polynomial Fock sections, with factor \(1\) in this range since the
Witt index of \(Y_0\) is at most \(r/2\) and \(r>2\ell\).
\cite[Theorem~2 and Lemma~3.1, pp.~237, 246]{tr3yamana}
The theta coefficient selects the Gram fibre. On the Eisenstein
side the nonsingular coefficient unfolds through the open cell.
At \(s_0\), the local section at \(w_\ell n(b)\), where \(w_\ell\) is
the long Siegel Weyl element, is the Fourier transform of the
pushforward of \(\Phi_\nu(X)\,dX\) under \(X\mapsto X^*JX\), times
the local Weil constant in the action of \(w_\ell\). Applying
\(n(b)\) and evaluating that action at zero gives the integral with
phase \(\vartheta_\nu(\operatorname{tr}(bX^*JX))\).
Fourier inversion for factorizable \(\Phi=\bigotimes_\nu\Phi_\nu\)
therefore gives
\[
  E_{\Phi,S}(s_0,1)
  =\int_{[G_0]}\ \sum_{\substack{X\in Y_0^{\ell}\\ X^*JX=S}}
     \Phi(h^{-1}X)\,dh
     =\prod_\nu \alpha_\nu(S;\Phi_\nu),
\]
where \(\alpha_\nu(S;\Phi_\nu)\) is the value at \(S\) of the density
of the pushforward of \(\Phi_\nu(X)\,dX\) under \(X\mapsto X^*JX\),
using the self-dual measures for \(\vartheta_\nu\) on the tuple space
with the trace of the Hermitian form and on the Hermitian matrices
with the trace pairing.
The Weil indices have global product \(1\).
The reductions of \(S\) can have different coranks at different finite
places. The formula only needs \(S\) to be nonsingular over \(E\), and
below it is applied only to Gram matrices of rational frames. In the
release the Siegel--Weil formula enters the scalar main term mainly
through the equality of compatible stabilizer volumes derived in the
paragraph \emph{Stabilizer volumes} below. The release also uses it
for the finiteness and the infinite--finite splitting of the orbital
measures in the normalizing constants
\releasecite{Formalism/LocalToGlobal/SiegelLatticeTheory/CokernelGLCounts/RankPointCounts/DualFamilyCounting/NcSSplitSk.lean}{37}{orbital measure splitting}. The scalar
pairs are not themselves rewritten as Eisenstein coefficients. Their
main term comes from the Poisson-summation bound in the next paragraph.

In this normalization the signed sum over the scalar pairs is bounded
by a sum of squared archimedean amplitudes. Here the real orbital
calculation
\releasecite{Development/AbelianVarietyTheory/Proofs/FlagMeasures/OffK0Closed.lean}{654}{uses the chosen local Bessel integrals of both words}.
Positivity and Poisson summation
\releasecite{Formalism/HermitianCurvature/HermitianCurvatureProofs/GrowthOfBasic.lean}{1136}{bound that sum}
through a recursion over the Poisson levels
\releasecite{Formalism/HermitianCurvature/StatementTB.lean}{160}{one Poisson step per level}
by \((1+\delta_{\mathrm{sc}})E_{\mathrm{den}}\) times
the archimedean energy integral, with \(0\le\delta_{\mathrm{sc}}\le b^{-7}\).
The inequality applies to the \(\Gamma_1\)-weighted ordered pairs in
the finite convolution. After disintegration the formal degrees bound
the archimedean integral by \(1/\mathfrak d\). Write \(V_{wv}\) for the
adelic volume of the joint stabilizer in the mass calculation,
\(V_w,V_v\) for the two adelic Bessel stabilizer volumes, and \(s_w,s_v\)
for their deep Bessel level volumes. Put
\(c_w=c_0(w)/v_{\rm lev}\), removing the common ambient level volume
\(v_{\rm lev}\) already supplied by the unfolded level integral.
The product of the local measures, including the finite factor
\(\Gamma_1\) at the places above the auxiliary primes,
\releasecite{Formalism/HermitianCurvature/HermitianCurvatureProofs/LocalVolumePositivity.lean}{95}{gives}
\(v^\dagger/(s_ws_v)\). The resulting bound for the \(\Gamma_1\)-weighted
signed sum is \((1+\delta_{\mathrm{sc}})E_{\mathrm{den}}\) times
\(V_{wv}v^\dagger/(s_ws_vc_wc_v\mathfrak d)\). The stabilizer volumes
and numbers of columns satisfy, as derived in the paragraphs
\emph{Stabilizer volumes} and \emph{Columns} below,
\[
  V_{wv}V_G=V_wV_v,\qquad
  c_w^{-1}=\frac{d_ws_w}{V_w},\qquad
  c_v^{-1}=\frac{d_vs_v}{V_v}.
\]
Consequently the scale multiplying the allowance is
\[
  \frac{V_{wv}v^\dagger}{s_ws_v c_wc_v\mathfrak d}
   =\frac{d_wd_v v^\dagger}{\mathfrak d V_G}=D_{wv}.
\]
The signed sum over the generic scalar pairs is therefore at most
\(A D_{wv}\), where \(A=(1+\delta_{\mathrm{sc}})E_{\mathrm{den}}\).
The release states the same bound with the factor \(P_g\) of the finite
convolution carried along
\releasecite{Formalism/HermitianCurvature/HermitianCurvatureProofs/LocalVolumePositivity.lean}{176}{mass formula with constant one}.

\leavevmode\teamanchor{proof.stabilizer-volume-identity}%
\emph{Stabilizer volumes.} Fix a Haar measure \(dg\) on
\(G_0(\mathbf A_F)\) and put \(V_G=\operatorname{vol}([G_0],dg)\).
Let \(a\ge1\) with \(r>2a\), let \(X_0\in Y_0^a\) have nonsingular Gram
matrix \(S=X_0^*JX_0\), and let \(G_{X_0}\subset G_0\) be its pointwise
stabilizer, the unitary group of the orthogonal complement of the span
of \(X_0\). Let \(\nu_S=\bigotimes_\nu\nu_{S,\nu}\) be the measure on the
adelic fibre \(\{X:X^*JX=S\}\) given by the local Gelfand--Leray
measures of the self-dual measures above, so that
\(\int\Phi\,d\nu_S=\prod_\nu\alpha_\nu(S;\Phi_\nu)\) for factorizable
\(\Phi\). Give \(G_{X_0}(\mathbf A_F)\) the Haar measure \(dh'\)
compatible with \(dg\) and \(\nu_S\), that is \(dg=dh'\,d\nu_S\) under
the orbit map \(g\mapsto gX_0\). The adelic stabilizer volumes of the
mass calculation are taken with
\releasecite{Formalism/PeriodIntegrals/VolNzeroOne.lean}{60}{this compatible normalization}.
We show that \(\operatorname{vol}([G_{X_0}],dh')=V_G\).

By Witt's theorem for Hermitian spaces over \(E\), the isometry between
the spans of two tuples with Gram matrix \(S\) that matches the tuples
extends to \(Y_0\). Hence the rational fibre is the single orbit
\(G_0(F)X_0\), with stabilizer \(G_{X_0}(F)\), and each local fibre is
one \(G_0(F_\nu)\)-orbit. At almost all finite places \(\nu\), the tuple
\(X_0\) lies in a self-dual lattice \(\Lambda_\nu\), \(S\) is unimodular,
the residue characteristic is odd and \(E_\nu/F_\nu\) is unramified or
split. There Witt's theorem over the residue field \(k_\nu\) and Hensel's
lemma show that the integral points of the fibre form one orbit of the
stabilizer of \(\Lambda_\nu\), because the Gram map is smooth at tuples
with unimodular Gram matrix. Thus the adelic fibre is
\(G_0(\mathbf A_F)X_0\).

Let \(\Phi=\bigotimes_\nu\Phi_\nu\ge0\) be a Schwartz function with
\(\Phi_\infty\) smooth and compactly supported. All terms are
nonnegative, so the integrals below may be rearranged. Replace \(g\) by
\(g^{-1}\) and integrate first over \(G_{X_0}(\mathbf A_F)/G_{X_0}(F)\).
This gives
\[
 \int_{[G_0]}\sum_{X^*JX=S}\Phi(g^{-1}X)\,dg
 =\int_{G_{X_0}(F)\backslash G_0(\mathbf A_F)}\Phi(g^{-1}X_0)\,dg
 =\operatorname{vol}([G_{X_0}],dh')\int\Phi\,d\nu_S .
\]
Dividing by \(V_G\) gives the theta side of the coefficient formula,
since that formula uses the measure of mass \(1\). The formula, with
\(\ell\) replaced by \(a\), is valid because \(r>2a\). It identifies
the quotient with \(\prod_\nu\alpha_\nu(S;\Phi_\nu)=\int\Phi\,d\nu_S\).
Thus
\(\operatorname{vol}([G_{X_0}],dh')\int\Phi\,d\nu_S=V_G\int\Phi\,d\nu_S\).

It remains to choose \(\Phi\) with \(0<\int\Phi\,d\nu_S<\infty\). Take
\(\Phi_\nu(X_0)>0\) at every place, and take \(\Phi_\nu\) to be the
indicator of \(\Lambda_\nu^a\) at the almost all places above. Since
\(X_0\) is injective, the Gram map is a submersion at \(X_0\). The
local fibre is therefore a manifold through \(X_0\) (analytic at finite
places) with a positive Gelfand--Leray measure, so
\(\alpha_\nu(S;\Phi_\nu)>0\). At the almost all places, the self-dual
measures give \(\Lambda_\nu^a\) and the integral Hermitian matrices
volume \(1\). Let \(q_\nu=\#k_\nu\). By Hensel's lemma
\(\alpha_\nu=\#\mathcal F(k_\nu)/q_\nu^{\,r^2-(r-a)^2}\), where
\(\mathcal F(k_\nu)\) is the fibre over \(k_\nu\). By Witt's theorem
over \(k_\nu\), this fibre is \(U_r(k_\nu)/U_{r-a}(k_\nu)\) in the
inert case and \(\mathrm{GL}_r(k_\nu)/\mathrm{GL}_{r-a}(k_\nu)\) in the
split case. The group orders
\(q^{n^2}\prod_{j=1}^n(1-(-q)^{-j})\) and
\(q^{n^2}\prod_{j=1}^n(1-q^{-j})\) then give
\[
 \alpha_\nu(S;\Phi_\nu)=\prod_{j=r-a+1}^{r}\bigl(1-\epsilon_\nu^{\,j}q_\nu^{-j}\bigr),
 \qquad \epsilon_\nu=\mp1 \text{ for inert, split } \nu .
\]
Since \(r-a+1\ge a+2\ge3\), the product over \(\nu\) converges
absolutely to a nonzero number. Hence
\(\operatorname{vol}([G_{X_0}],dh')=V_G\).

The Bessel stabilizer of \(w\) is the pointwise stabilizer of the
\((m-1)\)-tuple made of the isotropic frame, its dual frame and the
vector of norm \(-T_w\). Its Gram matrix is nonsingular, since the
frames pair to the identity and \(T_w\neq0\). A generic pair
has a nondegenerate joint span, and its joint stabilizer is the
pointwise stabilizer of an \(m\)-tuple. Here
\(r=k+m-1>2m\), because \(k\ge\ell+f+5\) and \(\ell=m\). Hence
\(V_w=V_v=V_{wv}=V_G\), which gives \(V_{wv}V_G=V_wV_v\). The release
records the identification of the mass-calculation stabilizers with
these frame stabilizers at
\releasecite{Formalism/LocalToGlobal/SiegelLatticeTheory/CokernelGLCounts/MassDataInst.lean}{257}{the volume input of the mass data}.

\leavevmode\teamanchor{proof.column-volume-identity}%
\emph{Columns.} Let \(N\) be the unipotent radical of the Bessel datum
of \(w\), with the self-dual measure, so that
\(\operatorname{vol}(N(F)\backslash N(\mathbf A_F))=1\). Let
\(v_N^\infty\) be the volume of a real fundamental domain for
\(N(F)\cap N(\mathbf R)(K_f\cap N)\), where \(K_f\) is the ambient
level. Then \(v_N(w)=v_N^\infty v_{\rm lev}\), so
\(c_w=v_N^\infty v_H(w)/\dim\pi_w\). Let \(K''_w\subset H_w(\mathbf A_f)\)
be the deep Bessel level. Write the compatible measure of the previous
paragraph on \(H_w(\mathbf A_F)\) as the product of a real measure and a
Haar measure \(\mu_H\) on \(H_w(\mathbf A_f)\). The release
\releasecite{Formalism/LocalToGlobal/SiegelLatticeTheory/CokernelGLCounts/RankPointCounts/MassDataXMasses.lean}{9}{ties these measures}
to those of \eqref{eq:period-gram}, so the real factor gives
\(H_w(\mathbf R)\) volume \(v_H(w)\) and \(N\) the measure defining
\(v_N(w)\). The deep Bessel level volume is
\(s_w=\mu_N(K_f\cap N)\,\mu_H(K''_w)\), with the factorization of the
Bessel level
\releasecite{Formalism/PeriodIntegrals/OrbitalBoundsConclusion/MassDataLevelS.lean}{5}{supplied by the release}.

First, \(N\) is unipotent, so strong approximation gives
\(N(\mathbf A_F)=N(F)N(\mathbf R)(K_f\cap N)\). Thus
\(1=v_N^\infty\mu_N(K_f\cap N)\).

Second, write
\(H_w(\mathbf A_f)=\bigsqcup_{c}H_w(F)\,c\,K''_w\) over \(h_w\) classes.
Each stabilizer \(\Gamma_c=H_w(F)\cap cK''_wc^{-1}\) is finite because
\(H_w(\mathbf R)\) is compact. That \(\Gamma_c=1\) is taken from the release
\releasecite{Formalism/LocalToGlobal/SiegelLatticeTheory/HaarTransversals/UnitaryHaarProducts/InHHGamUnipKu.lean}{12}{neatness of the deep level}.
It is the analogue of the pro-\(g_i\) argument used above for the
classes of level \(K_0\): if the deep level reduces into the upper
unitriangular group above each auxiliary prime, every element of
\(\Gamma_c\) has order a power of \(g_0\) and a power of \(g_1\). Hence
\(V_w=h_w\,v_H(w)\,\mu_H(K''_w)\).

Third, \(H_w^{\mathrm{full}}\) is the space of \(\pi_w\)-valued functions
on these classes
\releasecite{Formalism/LocalToGlobal/SiegelLatticeTheory/CokernelGLCounts/RankPointCounts/MassDataXColumns.lean}{18}{with the measures of the period Gram identity},
so \(d_w=h_w\dim\pi_w\).

Combining the three facts,
\[
 \frac{d_ws_w}{V_w}
 =\frac{\dim\pi_w\,\mu_N(K_f\cap N)}{v_H(w)}
 =\frac{\dim\pi_w}{v_N^\infty v_H(w)}=c_w^{-1}.
\]
The same argument applies to \(v\).

\leavevmode\teamanchor{proof.orbital-error-estimates}%
The absolute sums over pairs of corank at least two, and separately
over generic nonscalar pairs of corank one, are at most
\(\epsilon_L D_{wv}\) in each case, where
\(\epsilon_L=\exp(-kf\log k/3)\). The first sum includes both
nondegenerate families. These and the signed bound below are bounds
for the \(\Gamma_1\)-weighted ordered pairs, with the normalization of
the scalar term, so the factor \(P_g\Gamma_1\) of the finite
convolution becomes \(P_g\). Let \(\mathcal P\) be the union of the two
sets of pairs, and put \(x=kf\log k\). Each term of the absolute sums is
a pairing of two orbital vectors. By the Cauchy--Schwarz inequality it
is at most the product of their norms. The edition takes two inputs
from the release, both under the standing parameter ranges of the
family. First, decay of the real orbital amplitudes
\releasecite{Formalism/LocalToGlobal/SiegelLatticeTheory/FlagStabilizerTensors/RootLatticeThresholds/LoopsEv.lean}{18}{bounds the weighted sum of these products over \(\mathcal P\)}
by \(A_{\mathrm L}(\Sigma_{\rm nd}+\Sigma_{\rm deg}+\Sigma_1)D_{wv}\).
Here \(A_{\mathrm L}\ge1\) is a constant with
\(6\log(3A_{\mathrm L})\le x\), \(\Sigma_{\rm nd}\) and
\(\Sigma_{\rm deg}\) are the weighted lattice sums over the
nondegenerate and degenerate rows of coranks \(2,\ldots,k\) (in the
release the degenerate sum carries the square of a weight parameter,
which its lattice bound cancels), and
\(\Sigma_1\) is the contribution of a single row. Second, counting
lattices at fixed corank gives
\releasecite{Formalism/FlagUnipotentLattices/MinkowskiFactorialBounds/ExponentInequalityNondegenerateMid.lean}{609}{\(\Sigma_{\rm nd}\le e^{-4x}\)},
after the level comparison
\releasecite{Formalism/FlagUnipotentLattices/MinkowskiFactorialBounds/LatticeSumBounds.lean}{429}{for the nondegenerate rows},
\releasecite{Formalism/FlagUnipotentLattices/MinkowskiFactorialBounds/ExponentInequalityNondegenerateMid.lean}{671}{\(\Sigma_{\rm deg}\le e^{-x/2}\)}
and
\releasecite{Formalism/FlagUnipotentLattices/MinkowskiFactorialBounds/ExponentInequalityNondegenerateMid.lean}{735}{\(\Sigma_1\le e^{-x}\)}.
Each of the three bounds is at most \(e^{-x/2}\), so the weighted
absolute sum over \(\mathcal P\) is at most
\(3A_{\mathrm L}e^{-x/2}D_{wv}\). Since
\(\log(3A_{\mathrm L})\le x/6=x/2-x/3\), this is at most
\(e^{-x/3}D_{wv}=\epsilon_LD_{wv}\). The terms are nonnegative, so the
same bound holds for each of the two sets of pairs.

For nongeneric pairs of corank one, a failure of local generation
decreases either the mass of the local configurations or the
character sum. Summing the weights over the fibres of the local data
\releasecite{Formalism/LocalToGlobal/SiegelLatticeTheory/CokernelGLCounts/RowsOther.lean}{166}{bounds the signed sum}
by \(\epsilon_N D_{wv}\), where
\[
 \epsilon_N=2^{-(k-2r_S-5)}\exp(C_Nm^5 f\log b)
\]
for a fixed \(C_N\ge0\). The edition takes this signed bound from the
release. Its inputs are the finite-place data of the nongeneric
pairs (classes, amplitudes, normals, points, local generation, mass
and heights)
\releasecite{Formalism/LocalToGlobal/SiegelLatticeTheory/CokernelGLCounts/RowsOther.lean}{132}{finite-place inputs}
and the real measure and point-sum bounds
\releasecite{Formalism/LocalToGlobal/SiegelLatticeTheory/CokernelGLCounts/RowsOther.lean}{136}{real inputs}.

\leavevmode\teamanchor{proof.mixed-word-estimate}%
The generic row sum contains the scalar main term, the higher-corank
error and the generic nonscalar corank-one error. Its associated
operator has squared norm at most \(P_g(A+2\epsilon_L)D_{wv}\).
The nongeneric part contains the higher-corank and signed corank-one
errors, so its operator has squared norm at most
\(P_g(\epsilon_L+\epsilon_N)D_{wv}\). Hence
\[
 \|K_{wv}\|_{\mathrm{HS}}^2
 \le D_{wv}\left[
   \sqrt{P_g(A+2\epsilon_L)}+\sqrt{P_g(\epsilon_L+\epsilon_N)}
 \right]^2.
\]
The selected family has \(m^4\le f\le k\le x^5\),
\(T_w\ge b^{m/2}\), \(b\ge16\), \(g_i\ge(kb)^8\) for \(i=0,1\), and
\(k\ge\max(32,4r_S+12,70\log b)\). Each \(g_i\) has \(f\) places of
\(F\) above it, all of residue size \(g_i\), so
\[
 \log P_g\le 3f(g_0^{-1}+g_1^{-1})\le 6b^{-7}.
\]
Since \(C_N\) is fixed, the
\hyperlink{proof.host-realization}{size bounds of the family} give
\(8C_Nm^5 f\log b\le k\) beyond a threshold depending on \(C_N\).
Since \(k-2r_S-2\ge k/2+4\) and \(k-2r_S-5\ge k/2+1\),
the exponents in the density and error estimates give
\(\log E_{\mathrm{den}},\epsilon_L,\epsilon_N\le e^{-k/5}\le b^{-14}\).
In particular \(E_{\mathrm{den}}\le\exp(b^{-14})\).
Substituting these bounds in the two square roots
\releasecite{Formalism/HermitianCurvature/HermitianCurvatureProofs/GrowthOfBasic.lean}{232}{gives, for \(b\ge16\),}
\begin{equation}
 \label{eq:trace-mixed-block}
 \|K_{wv}\|_{\mathrm{HS}}^2
 \le (1+b^{-6})\,\frac{d_wd_v}{D_\star}\qquad(w\ne v).
\end{equation}

\subsection{Trace moments after restricting the columns}

\noindent\leavevmode\teamanchor{candidate.count}%
\textbf{AI-proposed candidate (uniform multiplicities for unitary groups at finite level).}
Let \(G=U(V)\) for a rank-\(r\) Hermitian space over a CM extension
\(E/F\) unramified at every finite place, with a self-dual lattice
\(\Lambda\). Put \(f=[F:\mathbb Q]\) and
\(\mathrm{rd}(E)=|\operatorname{disc}E|^{1/(2f)}\). The real signatures
are all \((p,q)\), with \(p+q=r\) and \(q^2\le r\). Let \(K_f\) be the
finite adelic stabilizer of \(\Lambda\) and \(K_f(g)\) its principal
congruence subgroup modulo \(g\). With the given finite Haar measure,
take a locally constant self-adjoint convolution idempotent \(e\)
of positive type, supported on \(K_f\) and bi-invariant under \(K_f(g)\).
For a rational prime \(g_0>4r\) splitting completely in \(E\), require
\((x-1)^r\widehat\Lambda\subseteq g_0\widehat\Lambda\) for \(x\) in
its support, where \(\widehat\Lambda\) is the finite adelic lattice.
These are the
\releasecite{Formalism/SharpRestriction/CuspSpanNormBounds/GlobAut.lean}{213}{finite projector conditions}.

For fixed \(A>0\) and \(r\) beyond a threshold depending on \(A\),
assume \(b\ge2\), \(1\le g\le b^A\), \(\log\log\mathrm{rd}(E)\le r\),
\(\log f\le\sqrt r\), and \(\log\log b\le rf\log r/16\).
At each real place \(v\), let \(\sigma_v\) be a discrete series with
strictly decreasing Harish--Chandra parameters
\(\lambda_{v,i}\in\mathbb Z+(r-1)/2\) and their signature labels.
Require \(|\lambda_{v,i}|\le b^2\) and, for every pair with opposite
signature labels,
\[
 |\lambda_{v,i}-\lambda_{v,j}|
 \ge 8\cdot10^{14}(A+4)\,r^{19} f^5\log(b\,\mathrm{rd}(E)).
\]
Write \(\sigma_\infty=\widehat\bigotimes_{v\mid\infty}\sigma_v\)
and \(R(e)\) for finite convolution in the right regular
representation on \(L^2([G])\). The count \(D(e,\sigma_\infty)\) is
the dimension of \(R(e)\operatorname{Hom}_{G_\infty}^{\rm bdd}
(\sigma_\infty,L^2([G]))\), which is finite. The estimate is
\[
 \begin{aligned}
 D_\star(e,\sigma_\infty)
    &=\operatorname{vol}([G])\,\operatorname{Re}e(1)
       \prod_{v\mid\infty}d(\sigma_v),\\
 |D(e,\sigma_\infty)-D_\star(e,\sigma_\infty)|
    &\le \exp(-rf\log r/16)\,D_\star(e,\sigma_\infty).
 \end{aligned}
\]
The formal degrees use the real Haar measures in the quotient volume.
The error is uniform in the fields, ranks, signatures and projectors
satisfying these hypotheses.

\leavevmode\teamanchor{proof.count-residual}%
\paragraph{What the edition takes from the release.}
The edition takes this estimate, including the finiteness of
\(D(e,\sigma_\infty)\), from the release's
\releasecite{Development/AbelianVarietyTheory/Proofs/FlagMeasures/UnitaryAutomorphicCountingUnbundled.lean}{40}{counting statement}.
That statement repackages the release's
\releasecite{Development/AbelianVarietyTheory/BlockTheory/SiegelWeilLocal/LatticeCountingTheorem.lean}{10}{counting theorem},
with the same main term and error. The release proves the counting
theorem from a
\releasecite{Development/AbelianVarietyTheory/Statements/CuspFrameWeights/SameKindIsometric.lean}{32}{volume step}
together with
\releasecite{Development/AbelianVarietyTheory/Proofs/FlagMeasures/UnitaryAutomorphicCountingOfCompactDualVolumeRat.lean}{37}{kernel bounds},
and passes to the stated form through its
\releasecite{Development/AbelianVarietyTheory/Proofs/FlagMeasures/CountTypedOfVolumeStep.lean}{13}{lattice-counting input from the volume and kernel estimates}.
The volume step compares the normalized adelic volumes of \(U(V)\) and
of \(U(x^\perp)\) for a vector \(x\) with
\(\langle x,x\rangle=\pm1\) and \(\operatorname{rank}V\ge3\). The
release derives it from the convergent unitary Siegel--Weil formula
\releasecite{Development/AbelianVarietyTheory/Proofs/FlagMeasures/VolumeStepOfSiegelWeil.lean}{13}{volume step from Siegel--Weil}.
The edition has not compared that input with the published convergent
case \cite{tr3yamana}.

The shape of the main term is explained as follows. Put
\(\phi=e\otimes q_\infty\), where \(q_\infty\) is the real projector
of the row realization below. By that realization,
\(D(e,\sigma_\infty)\) is the rank of the orthogonal projection
\(R(\phi)\) on \(L^2([G])\). Its kernel is
\(\sum_{\gamma\in G(F)}\phi(x^{-1}\gamma y)\) wherever this series
converges. Its identity term on the diagonal integrates to
\(\operatorname{vol}([G])\,e(1)\,q_\infty(1)=D_\star(e,\sigma_\infty)\),
because \(q_\infty(1)=\prod_v d(\sigma_v)\) and \(e(1)\) is real for
self-adjoint \(e\). The release controls the convergence and the
non-identity terms. It uses large-gap decay of lowest-type matrix
coefficients, isotropic and nonzero flag probability bounds derived
from the volume step, near-unipotent and far-term bounds, and a
kernel bound on the deep part of the quotient. The edition has not
reconstructed these bounds. The residual implication is exactly the
displayed estimate under the stated hypotheses.

Shin--Templier's quantitative Plancherel results and counting
asymptotic apply to level and weight families of a fixed group over
a fixed base field. \cite[Theorems~9.16 and~9.19 and Corollary~9.25]{tr3shin-templier}
The bound here is uniform as the unitary rank and field vary,
and its large-gap condition permits small gaps within a signature
block. It concerns the full projected multiplicity at the given
real representation.
In the application below the quotient \([G]\) is not compact, so the
compact-quotient trace formula does not apply directly either.
There \(q=m/2\ge2\), so \(V\) is isotropic at every real place.
The vectors \(x\) with \(\langle x,x\rangle=0\) are the zeros of the
nondegenerate quadratic form \(x\mapsto\langle x,x\rangle\) on \(V\)
regarded as an \(F\)-space of dimension \(2r\). At every finite place
\(v\) of \(F\) this form is isotropic over \(F_v\), being a quadratic
form in \(2r\ge6>4\) variables over a nonarchimedean local field. At
the real places it has signature \((2p,2q)\) with \(p,q\ge1\). The
Hasse--Minkowski theorem therefore makes \(V\) isotropic over \(E\). Hence
\(G\) has positive \(F\)-rank and \([G]\) is noncompact.

\leavevmode\teamanchor{proof.full-row-count}%
We apply the multiplicity input to the full row projector of the
same record as \(\mathcal P_w\). The finite quotient at level
\(g_0g_1\) of the selected self-dual lattice stabilizer is
\(\prod_{i=0}^1\prod_{z\mid g_i}\mathrm{GL}_r(\mathbf F_{g_i})\)
by the split reductions and the Chinese remainder theorem.
The full row projector \(e\) is the tensor product of the finite
character projectors on the prescribed unitriangular subgroups at
these places, inflated with the principal congruence kernel of
measure \(\operatorname{vol}K_f(g_0g_1)\).
Equivalently \(e\) is the average with conjugate character over the
deep ambient subgroup of volume \(v^\dagger\). Its identity value is
\(1/v^\dagger\). With this reciprocal-volume normalization it is
self-adjoint and idempotent, and it is of positive type as the
matrix coefficient of the vector \(e\) in the regular
representation. Its support reduces to a unipotent group at each
auxiliary prime, so \((y-1)^r\widehat{\Lambda}\subseteq
g_i\widehat{\Lambda}\) for every point \(y\) of the support.
The condition also holds on the conjugate half of each split
unitary reduction, which is the inverse transpose of the chosen half.

\leavevmode\teamanchor{proof.row-realization}%
At the real places, evaluation of the bounded intertwiner at the
chosen tensor of unit lowest-type vectors identifies the
projected multiplicity space with the Petersson row model
\(\mathcal H\) in \eqref{eq:period-gram}. We give the argument.
Let \(\pi_\infty\) be the actual ordered tensor of real discrete
series on the Hilbert tensor product \(H_\infty\), and \(u\) its
chosen unit lowest-type tensor. It is irreducible, so \(u\) is
cyclic. Let \(d(\pi_\infty)=\prod_vd(\sigma_v)\) be its formal degree
for the real Haar measure used in the quotient volume. The adelic
group is unimodular, so right translation \(R\) is unitary on
\(L^2([G])\). For an integrable \(\varphi\) on \(G_\infty\), or a
locally constant compactly supported \(\varphi\) on the finite adelic
group, put \(R(\varphi)=\int\varphi(g)R(g)\,dg\). Then
\(\|R(\varphi)\|\le\|\varphi\|_1\), \(R(\varphi)^*=R(\varphi^*)\) with
\(\varphi^*(g)=\overline{\varphi(g^{-1})}\), and
\(R(\varphi_1)R(\varphi_2)=R(\varphi_1*\varphi_2)\) by Fubini.
Real and finite convolutions commute. In the right-translation
convention use
\[
 q_\infty(g)=d(\pi_\infty)
                  \overline{\langle u,\pi_\infty(g)u\rangle}.
\]
The matrix coefficient is integrable at the chosen real gaps. Indeed
it is a product over the real places of coefficients of the
\(K\)-finite lowest-type vectors of discrete series of \(U(p,q)\).
Since \(U(p,q)\) is the product of its compact center and
\(SU(p,q)\), and the center acts by scalars of modulus one, it
suffices to integrate over \(SU(p,q)\). This is a connected
semisimple group with finite center, with the same roots, and the
restriction is a discrete series with the same Harish--Chandra
parameter. For a noncompact root \(\alpha=e_i-e_j\) in the
normalization with roots of squared length two, the sum over all
roots gives \(\tfrac14\sum_\beta|(\alpha,\beta)|=r-1\). The gaps
\(|\lambda_{v,i}-\lambda_{v,j}|\) between oppositely labelled
parameters exceed \(r-1\). By the size checks below they are at least
\(C_{\rm gap}r^{19}f^5L\), and \(L=\log b\ge\log2\). Hence the
\(L^1\) criterion for discrete-series coefficients
\cite[Part~III, Corollary to Theorem~1]{p2milicic-integrability}
applies at each real place. Tonelli's theorem gives integrability of
the product.

Schur orthogonality for the square-integrable irreducible
representation \(\pi_\infty\), with the inner product linear in its
second variable, reads
\[
 \int_{G_\infty}\langle b,\pi_\infty(g)a\rangle
   \overline{\langle e',\pi_\infty(g)c\rangle}\,dg
   =d(\pi_\infty)^{-1}\langle c,a\rangle\langle b,e'\rangle .
\]
Unitarity gives \(q_\infty^*(g)=d(\pi_\infty)\langle
\pi_\infty(g)u,u\rangle=q_\infty(g)\). For \(g,x\in G_\infty\),
\[
 (q_\infty*\delta_g*q_\infty)(x)
 =d(\pi_\infty)^2\int_{G_\infty}
   \overline{\langle u,\pi_\infty(h)u\rangle}\,
   \langle \pi_\infty(x)u,\pi_\infty(h)\pi_\infty(g)u\rangle\,dh .
\]
Schur orthogonality with \(a=b=u\), \(c=\pi_\infty(g)u\) and
\(e'=\pi_\infty(x)u\), followed by complex conjugation, evaluates the
integral as
\(d(\pi_\infty)^{-1}\langle u,\pi_\infty(g)u\rangle
\overline{\langle u,\pi_\infty(x)u\rangle}\). Hence
\[
 q_\infty*\delta_g*q_\infty
                =\langle u,\pi_\infty(g)u\rangle q_\infty ,
\]
and \(g=1\) gives \(q_\infty*q_\infty=q_\infty\). Thus
\(T=R(q_\infty)\) is an orthogonal projection, commuting with
\(R(e)\), and \(TR(g)T=\langle u,\pi_\infty(g)u\rangle T\).
Also \(\int q_\infty(g)\pi_\infty(g)u\,dg=u\): pair with any \(y\)
and apply Schur orthogonality with \(b=y\) and \(a=c=e'=u\).
The row model \(\mathcal H\) is the common range of \(R(e)\) and
\(T\) at the fixed bottom level and real type.

Let \(A:H_\infty\to L^2([G])\) be a bounded intertwiner. Then
\(R(e)A\) is again an intertwiner, and the projected multiplicity
space consists of the intertwiners with \(R(e)A=A\). For such \(A\),
\(TAu=A\int q_\infty(g)\pi_\infty(g)u\,dg=Au\) and \(R(e)Au=Au\), so
\(Au\in\mathcal H\). If \(Au=0\), then \(A\) vanishes on the orbit of
\(u\) and hence on \(H_\infty\) by cyclicity. Conversely let
\(z\in\mathcal H\). For \(a,b\in G_\infty\),
\[
 \langle R(a)z,R(b)z\rangle
 =\langle z,TR(a^{-1}b)Tz\rangle
 =\|z\|^2\langle\pi_\infty(a)u,\pi_\infty(b)u\rangle .
\]
So the Gram of the image of any finite orbit combination is its
original Gram multiplied by \(\|z\|^2\). The prescription
\(\pi_\infty(g)u\mapsto R(g)z\) is therefore well defined on the
span of the orbit, bounded by \(\|z\|\), and extends uniquely to a
bounded map \(A\) on \(H_\infty\) by density. It intertwines on the
orbit, hence everywhere by continuity. It satisfies \(Au=z\) and
\(R(e)A=A\), since \(R(e)R(g)z=R(g)R(e)z=R(g)z\) for real \(g\).
Evaluation at \(u\) is thus a linear bijection from the projected
multiplicity space onto \(\mathcal H\), with the Petersson norm of
\(z\) for the same Haar. The dimension \(D\) is therefore the
multiplicity of exactly the full row projector to which the count
applies.

Here are the size checks for this application and for the restricted
projectors below. The real parameters, in order, are a decreasing
nonnegative block \(v_1,\ldots,v_{k-1}\), all with the positive label,
and \(-B_1,\ldots,-B_m\) with alternating labels starting with the
positive label. Thus the signature is
\((k-1+m/2,m/2)\). The selected bound \(m^2\le k\)
gives \((m/2)^2\le r\), where \(r=k+m-1\).
Write \(L=\log b\) and \(\Delta\) for the weight displacement.
The selected inequalities for the weights in this real type give
\[
  0\le v_j\le4k\Delta,\qquad
  C_{\rm gap}mk\Delta\le B_1,\qquad
  B_m\le b^2,\qquad
  B_{j+1}-B_j\ge C_{\rm gap}r^{19}f^5 L .
\]
\releasecite{Formalism/SharpRestriction/CuspSpanNormBounds/StandingNumericalData.lean}{106}{Weight and level inequalities}.
\leavevmode\teamanchor{proof.count-parameters}%
The family comparison allows \(C_{\rm gap}\) at least
\(\max(2,16\cdot10^{14}(A+4))\), with \(A=2C_g\) when \(g_i\le b^{C_g}\).
Here \(C_g\), hence \(A\), is fixed on the family. We check each
hypothesis of the count at every real place.
\begin{enumerate}
\item The parameters are the Harish--Chandra parameters of the real
discrete series in the row type, so they lie in
\(\mathbb Z+(r-1)/2\) with the stated labels. They are regular,
hence distinct, and strictly decreasing in the stated order, since
\(B_1\ge C_{\rm gap}mk\Delta>0\).
\item Since \(C_{\rm gap}m\ge4\), we have \(v_j\le4k\Delta\le B_1\le
B_m\le b^2\). So all parameters have absolute value at most \(b^2\).
\item The labels of \(-B_i\) alternate starting with the positive
label, so \(-B_i\) has the negative label exactly for even \(i\).
Two oppositely labelled negative parameters \(-B_i,-B_j\) with
\(i<j\) satisfy \(B_j-B_i\ge B_{i+1}-B_i\), one adjacent \(B\)-gap.
A pair \(v_j,-B_i\) with opposite labels has \(i\) even, so
\(v_j+B_i\ge B_2\ge B_2-B_1\), again an adjacent gap. Pairs with
the same label impose no condition. Since
\(\mathrm{rd}(E)\ge1\) and \(\log\mathrm{rd}(E)\le L/C_{\rm gap}\le
L\), we have \(\log(b\,\mathrm{rd}(E))\le2L\). Hence every required
gap is at least
\(C_{\rm gap}r^{19}f^5L\ge 8\cdot10^{14}(A+4)\,r^{19}f^5\log(b\,
\mathrm{rd}(E))\).
\item The signature \((k-1+m/2,m/2)\) is the same at every real place,
and \(q=m/2\) satisfies \(q^2\le m^2\le k\le r\).
\item The selected levels are distinct rational primes, completely
split in \(E\), with \((kb)^8\le g_i\le b^{C_g}\). Hence
\(g_i\ge k^8>8k\ge4r\), \(1\le g_0g_1\le b^{2C_g}=b^A\), and
\(E/F\) is unramified above \(g_0\).
\item With \(b=\exp(x^\varepsilon)\), \(f\le x^5\) and
\(x^4\le k\le r\le2k\), we have \(b\ge2\),
\(\log\log b=\varepsilon\log x\le rf\log r/16\) and
\(\log f\le5\log x\le x^2\le\sqrt r\) for large \(x\). If
\(\log\mathrm{rd}(E)>1\), then
\(\log\log\mathrm{rd}(E)\le\log L=\varepsilon\log x\le r\); otherwise
\(\log\log\mathrm{rd}(E)\le0\). Finally \(r\ge x^4\) exceeds the rank
threshold for this fixed \(A\) once \(x\) is large.
\end{enumerate}
The CM extension is unramified at finite places and the selected
Hermitian lattice is self-dual, as in the ambient row construction.
The finite projector conditions were checked above. Thus all
hypotheses of the count hold for the common record at every
indicated translate for sufficiently large \(x\).

The main term is therefore \(D_\star=\mathfrak d V_G/v^\dagger\)
in the Gram normalization. The multiplicity input gives
\(|D-D_\star|\le\exp(-rf\log r/16)D_\star\), in particular \(D>0\).

\leavevmode\teamanchor{proof.full-trace-estimates}%
At each translate the estimates give
\begin{equation}
\label{eq:trace-full-gram}
 \begin{aligned}
  \|K_{wv}\|_{\rm HS}^2
       &\le (1+\eta_0)\frac{d_wd_v}{D_\star}
       && (w,v\in W,\ w\ne v),\\
  \|K_{ww}-I\|_{\rm HS}^2 &\le e_0 d_w &&(w\in W),\\
  \sum_{w\in W}\|K_{ww}-I\|_{\rm HS}^2 &\le e_0 R_{\rm full},
       &\qquad D&\le (1+\eta_1)D_\star .
 \end{aligned}
\end{equation}
Take \(\eta_0=b^{-6}\), \(\eta_1=\exp(-rf\log r/16)\),
and \(e_0\) from \eqref{eq:trace-single-diagonal}. The three
quantities \(\eta_0,\eta_1,\sqrt{e_0}\) are each bounded by
\(\exp(-x^{\varepsilon/8})\) at every indicated translate above a
common threshold.

\teamanchor{candidate.trace}%
\textbf{AI-proposed candidate (a uniform mixed Gram bound).}
For the fixed target, precision and tied family, the mixed estimate
\eqref{eq:trace-mixed-block} holds beyond one scale threshold at every
selected translate, simultaneously for all distinct words \(w,v\) in
its word set. It bounds \(\|L_w^*L_v\|_{\rm HS}^2\) by
\((1+b^{-6})d_wd_v/D_\star\) on the full columns, with the common
Petersson target and reference dimension at that translate.
For the selected seeds under the prime, weight-gap, lattice and rank
conditions above, the degenerate terms vanish. The upper bounds use
the signed sums for generic scalar pairs and nongeneric corank-one
pairs, and absolute sums for the remaining errors.

Getz--Hahn's compact-test formula supplies the orbit reindexing in a
simple relative trace formula. A quantitative precedent is
Kowalski--Saha--Tsimerman's estimate for Fourier coefficients of
genus-two Siegel Poincar\'e series over \(\mathbb Q\). For fixed
discriminant \(d\) their normalized error at indices \(L,M\) and
scalar weight \(\kappa\) is
\(O_{d,\epsilon}(L^{1+\epsilon}M^{3/2+\epsilon}\kappa^{-2/3})\).
\cite[Proposition~3.3, arXiv v4]{tr3kowalski-saha-tsimerman}
The unitary bound here uses the different local test and volume
identifications above, including the reference dimension from
\(V_G\), \(v^\dagger\) and \(\mathfrak d\). Its uniformity covers the
fields, ranks and word primes of the tied family. An aggregate trace
moment alone allows two identical isometric blocks of dimension
\(d\) inside a tight frame in dimension \(D_\star=2d\). Their mixed
squared norm is \(d\), greater than \((1+b^{-6})d/2\).

The trace argument sums the mixed bounds and uses the independent
single-word defect \eqref{eq:trace-single-diagonal} to control the
diagonal blocks. With the row and word counts these estimates give
the moments for Waldron's inequality. The ramified compression gives
\eqref{eq:trace-frame}, and \eqref{eq:trace-restricted-corank} controls
the rank after further column restrictions.

\leavevmode\teamanchor{proof.full-trace-identity}%
Write \(Q_{\rm full}=\sum_w L_w L_w^*\).
Cyclicity of trace gives
\[
 \operatorname{tr}(Q_{\rm full}^2)
 =\sum_{w,v}\operatorname{tr}(L_v^*L_w L_w^* L_v)
 =\sum_{w,v}\|K_{wv}\|_{\rm HS}^2.
\]
Put \(\delta_0=\sqrt{e_0}\). The diagonal terms have
\(\|K_{ww}\|_{\rm HS}\le (1+\delta_0)\sqrt{d_w}\) by the triangle
inequality. Summing the off-diagonal bound in \eqref{eq:trace-full-gram}
over ordered pairs \(w\ne v\) and using
\(\sum_{w\ne v}d_wd_v\le R_{\rm full}^2\) yields
\[
 \operatorname{tr}(Q_{\rm full}^2)
 \le (1+\eta_0)\frac{R_{\rm full}^2}{D_\star}
       +(1+\delta_0)^2 R_{\rm full}.
\]
\leavevmode\teamanchor{proof.trace-word-count}%
Section~\ref{sec:assembly} selects the word set with
\(|W|\ge C_0\Gamma^2 x^{C_H}\). These are the same words and
the same \(C_0=3\) as in the full column balance. This is part of the
\hyperlink{proof.choose-compatible-family}{simultaneous family selection}.
Together with the lower bound in
\eqref{eq:trace-balance} this gives
\(R_{\mathrm{full}}=\sum_wd_w\ge|W|D_\star/(C_0\Gamma^2)\ge x^{C_H}D_\star=D_\star/\sigma\),
where \(\sigma=x^{-C_H}\) and \(C_H>0\) is fixed on the chosen family.
This makes the diagonal term small compared with
\(R_{\mathrm{full}}^2/D_\star\).

\leavevmode\teamanchor{proof.compressed-frame-moments}%
For any isometric inclusions \(J_w:H_w^{\mathrm{act}}\hookrightarrow
H_w^{\mathrm{full}}\), define the restricted column maps \(A_w\) and
their frame operator \(Q_{\mathcal H}\) by
\[
 A_w=L_wJ_w,\qquad
 Q_{\mathcal H}=\sum_w A_w A_w^*
    =\sum_w c_0(w)^{-1}(\mathcal P_wJ_w)(\mathcal P_wJ_w)^* .
\]
Thus the frame operator weights the original column maps by
\(c_0(w)^{-1}>0\). The period comparison below uses \(J_w=\jmath_w U_w\).
Let \(R_{\mathrm{act}}=\sum_w r_w\) and
\(\rho=1-R_{\mathrm{act}}/R_{\mathrm{full}}\).
By \eqref{eq:trace-ramification}, \(0\le\rho\le2f/b\).

For every block, \(A_w^*A_v=J_w^*K_{wv}J_v\). In bases extending the
ranges of the \(J_w\), this compression retains a submatrix, and its
Hilbert--Schmidt norm is bounded by the full norm. Consequently the same
trace identity gives \(\operatorname{tr}(Q_{\mathcal H}^2)\le
\operatorname{tr}(Q_{\rm full}^2)\). Since \(J_w^*J_w=I\), compression
also bounds the sum of squared diagonal defects by \(e_0R_{\rm full}\).
Taking traces of the compressed diagonal blocks and applying
Hilbert--Schmidt Cauchy--Schwarz on their direct sum gives
\[
 t:=\operatorname{tr}Q_{\mathcal H}
 \ \ge\ R_{\rm act}-\sqrt{e_0R_{\rm act}R_{\rm full}}
 \ \ge\ (1-\rho-\delta_0)R_{\rm full}.
\]
For all sufficiently large \(x\) we have \(\sigma\le 1/20\),
\(\eta_0,\eta_1,\delta_0\le \sigma/2\) and \(\rho\le 1/4\), uniformly
over the indicated translates. In this range the last bound is
positive. Using \(D\le (1+\eta_1)D_\star\), we obtain
\[
 \frac{D\operatorname{tr}(Q_{\mathcal H}^2)}{t^2}
 \le
 \frac{(1+\eta_1)\bigl[1+\eta_0+(1+\delta_0)^2\sigma\bigr]}
      {(1-\rho-\delta_0)^2}.
\]
Write \(q_0=\sqrt{\rho}\). The numerator is at most
\((1+\sigma/2)^2+\sigma(1+\sigma/2)^3\le 1+3\sigma\), and the
denominator is at least \(1-\sigma-2q_0^2\). In the given ranges,
\[
 \begin{split}
 &(1+5\sigma+3q_0)(1-\sigma-2q_0^2)-(1+3\sigma)\\
 &\qquad =\sigma(1-5\sigma)
    +q_0\bigl(3(1-\sigma)-2q_0(1+5\sigma+3q_0)\bigr)\ \ge\ 0 .
 \end{split}
\]
Thus \(e=2e_0\) and \(\eta=5\sigma+3\sqrt{\rho}\) satisfy
\begin{equation}
\label{eq:trace-frame}
 \begin{gathered}
 \sum_w \|A_w^*A_w-I\|_{\rm HS}^2\le e R_{\rm act},
 \qquad 0\le e\le\tfrac14,\\
 \operatorname{tr}(Q_{\mathcal H}^2)
 \le (1+\eta)\frac{(\operatorname{tr}Q_{\mathcal H})^2}{D},
 \qquad \eta\ge0 .
 \end{gathered}
\end{equation}

The ramified rank bound \(\rho\le2f/b\), together with \(f\le x^5\)
and \(b=\exp(x^\varepsilon)\), gives
\(\sqrt\rho\le\exp(-x^{\varepsilon/8})\) for sufficiently large \(x\). With
\(E_x^0=\exp(-x^{\varepsilon/8})\), the bounds
\(\bar e=2(E_x^0)^2\) and \(\bar\eta=5x^{-C_H}+3E_x^0\) work
simultaneously at all indicated translates at scale \(x\).
Since \(C_H\) is fixed,
\begin{equation}
 \label{eq:trace-accuracy}
 (\eta+4\sqrt e)x^{C_H}
 \le(\bar\eta+4\sqrt{\bar e})x^{C_H}
 =5+(3+4\sqrt2)E_x^0x^{C_H}\le6
\end{equation}
once \(x\) is sufficiently large.

\subsection{The span of restricted columns}
\leavevmode\teamanchor{proof.analytic-span}%
For a finite sequence of complex Hilbert space vectors \(a_1,\ldots,a_N\)
with at least one nonzero vector, let \(d\) be the dimension of its span.
The generalized Welch inequality for vectors of arbitrary norms states
\[
 \left(\sum_{i=1}^N\|a_i\|^2\right)^2
 \le d\sum_{i,j=1}^N|\langle a_i,a_j\rangle|^2 .
\]
\cite[Theorem~1.3 and its proof, author version]{tr3waldron}
For the images of orthonormal bases under the \(A_w\), the established
positivity of \(t\) ensures a nonzero vector. The inequality gives
\(t^2\le d\operatorname{tr}(Q_{\mathcal H}^2)\) with
\(d=\operatorname{rank}Q_{\mathcal H}\).
By \eqref{eq:trace-frame}, \(d\ge D/(1+\eta)\).
With \(J_w=I\) the same argument gives \(D/(1+5\sigma)\).

For \(0\le\alpha\le1\), let \(S_w\subseteq H_w^{\mathrm{act}}\) be
complex subspaces with
\(\sum_w\operatorname{codim}S_w\le\alpha R_{\mathrm{act}}\).
Put \(M=\sum_w A_w S_w\), let \(P\) project orthogonally onto \(M^\perp\),
and let \(E_w\) project onto \(S_w^\perp\). Since \(P A_w=P A_w E_w\),
\[
 \begin{split}
 \operatorname{tr}(P Q_{\mathcal H})
 &=\sum_w\|P A_w E_w\|_{\rm HS}^2
   \le \sum_w \operatorname{tr}(E_w A_w^*A_w)\\
 &\le \sum_w\operatorname{rank} E_w+
   \sqrt{\sum_w\operatorname{rank} E_w}\,
   \sqrt{\sum_w\|A_w^*A_w-I\|_{\rm HS}^2}
   \le (\alpha+\sqrt e)R_{\rm act}.
 \end{split}
\]
From \eqref{eq:trace-frame}, \(t\ge (1-\sqrt e)R_{\rm act}\ge
R_{\rm act}/2\). Set \(z=\operatorname{tr}(P Q_{\mathcal H})/t\) and
\(p=\operatorname{rank}(P)/D\). Cauchy--Schwarz for the Hilbert--Schmidt
pairing of \(I-P\) with \(Q_{\mathcal H}\) gives
\((1-z)^2\le (1-p)(1+\eta)\). Since \(0\le z\le 1\),
\(p\le\eta+2z\). Therefore
\begin{equation}
\teamanchor{proof.restricted-corank}
\label{eq:trace-restricted-corank}
 \operatorname{codim}_{\mathcal H} M \le (\eta+4\alpha+4\sqrt e)D .
\end{equation}

\subsection{Projection to the common row dual}
\leavevmode\teamanchor{proof.rational-corank}%
We now use \(J_w=\jmath_wU_w\). Let \(\Pi:\mathcal H\to H_{\rm row}\)
be the orthogonal projection to \(H_{\rm row}=s(VE_{\rm row})\)
from \eqref{eq:period-block}. Since \(s\) is conjugate-linear and the
inner product is conjugate-linear in its first argument, the row
isometry gives a complex linear isomorphism
\(\mathcal T:H_{\rm row}\to E_{\rm row}^{\vee}\)
by \(\mathcal T(y)(F)=\langle s(VF),y\rangle_{\mathcal H}\).
The common \(V\) gives the same map \(\mathcal T\Pi\) for every
word. By \eqref{eq:period-bessel} the complexification of
\(B_{K,w}\) at \(\iota_K\) is
\[
 B_{K,w,\mathbf C}=a_w\sqrt{c_0(w)}\,\mathcal T\Pi A_w .
\]
In complexified row and column coordinates this is
\eqref{eq:period-block}, with the common row change and the
wordwise changes of basis specified there.
Taking the conjugate embedding at the same infinite place conjugates
all entries in the identity. Each nonzero scalar
\(a_w\sqrt{c_0(w)}\) preserves the complex image of the restricted map.

Write \(C_K=(R\otimes K)^\vee=C\otimes K\) and
\(D'=\dim_K C_K=\dim H_{\rm row}\).
For \(K\)-subspaces \(S_w\subset N_w\otimes K\), take their scalar
extensions \(S_{w,\iota_K}\subset H_w^{\mathrm{act}}\) in the prescribed
column models. Their complex codimensions are the original
\(K\)-codimensions. Put \(M=\sum_w A_wS_{w,\iota_K}\).
There is a surjection \(\mathcal H/M\to H_{\rm row}/\Pi M\).
If \(D'\ge(1-\delta)D>0\) for \(0\le\delta<1\) and
\(\sum_w\operatorname{codim}_K S_w\le\alpha R_{\rm act}\), then
\eqref{eq:trace-restricted-corank} and preservation of rank under
field extension give
\begin{equation}
 \label{eq:trace-rational-corank}
 \frac{\operatorname{codim}_{C_K}\sum_w B_{K,w}(S_w)}{D'}
 \le \frac{\eta+4\alpha+4\sqrt e}{1-\delta},
 \qquad 0\le\alpha\le1 .
\end{equation}
\subsection{Counting the distinguished rows}
\leavevmode\teamanchor{proof.distinguished-row-count}%
At each \(g_i\), let \(\boldsymbol{P}=(P_z)_{z\mid g_i}\) be the tuple of
local polynomial labels, with one factor for each of the \(f\) places
of \(F\) above \(g_i\). These label the finite-group types on the level
\(g_0g_1\) reduction of the ambient row space. For
\(g=g_i\), denote by \(z_P\in\mathbf C[\mathrm{GL}_r(\mathbf F_g)]\)
the central idempotent of the sum of irreducibles with label \(P\).
The labels are the monic degree-\(r\) polynomials with nonzero constant
term. The central projectors are orthogonal and sum to the identity.
For a set \(S\) of label tuples at \(g_i\) let \(e_{i,S}\) be the full
row projector \(e\) multiplied by the central projector
\(\sum_{\boldsymbol{P}\in S}\bigotimes_{z\mid g_i}z_{P_z}\) in the finite
quotient. All other local factors remain the full row factors. The
central projector commutes with \(e\). Thus multiplication selects
the allowed label spaces inside the range of \(e\). In particular
the resulting multiplicity counts the preimage of \(S\) in a basis
of full rows chosen in the joint label spaces. Repeated labels contribute
the dimension of their entire eigenspace. The finite projection in
the distinguished row construction is the central idempotent for
good labels at both primes. Its operator is the product of the
two single-prime good-label projectors, since on each joint
summand their actions are multiplication by the indicator functions
of the two good conditions. In the common Petersson model its
range is \(H_{\rm row}=s(VE_{\rm row})\).
\releasecite{Formalism/LocalToGlobal/SiegelLatticeTheory/FlagStabilizerTensors/RootLatticeThresholds/RowsGoodRed.lean}{39}{The common row model in the joint label spaces}.
\releasecite{Formalism/SharpRestriction/LevelLetterCocycles/ArithmeticCylinders.lean}{271}{Full and distinguished row indices}.

The joint good cut in this multiplicity space is cuspidal.
Let \(P=M_PU_P\) be a proper rational parabolic with first
block size \(0<a<r\). For each good tuple some split place
\(z\mid g_0\) has a label for which no subset of the
irreducible factor degrees sums to \(a\). For \(a>r/2\) use
the complement \(r-a\). The finite Green type at that place has
no invariants under the unipotent radical of the \((a,r-a)\)
parabolic. Indeed, Mackey's formula would put each cuspidal
block entirely on one side of the cut. The radical of
the actual local \(P_z\) contains this radical up to conjugacy
in the full compact stabilizer \(K_z\), since an integral basis
adapted to its flag makes the parabolic standard.

For a smooth vector in this finite isotypic space, write \(f_P\)
for its constant term. Fix all coordinates other than \(z\).
At an argument with \(z\)-coordinate \(p_z\in P_z\),
evaluation of \(f_P\) is invariant under right action on the
vector by \(U_{P,z}\cap K_z\). Indeed \(p_z\) normalizes
\(U_{P,z}\) and the constant term is left \(U_P(\mathbf A_F)\)-
invariant. Averaging the finite \(K_z\)-type over this
unipotent subgroup therefore gives zero for the evaluation.
Iwasawa \(G_z=P_zK_z\) gives vanishing for every \(z\)-coordinate.
The real vectors in question are smooth because they are
evaluations of bounded intertwiners on the smooth vector \(u\),
and the quotient in the unipotent integral is compact. Hence
\(f_P=0\) for every such \(P\). Summing over good tuples proves
the assertion on the exact common cut.

For the local character projector \(a_z\) of the full row space at
\(z\mid g_i\), put \(\omega_z(P)=(a_z z_P)[1]/a_z[1]\), where
\([1]\) takes the identity coefficient. In the regular representation
of a finite group \(G_f^{\rm red}\), the trace of left multiplication
by \(a\) is \(|G_f^{\rm red}|a[1]\). Thus \(\omega_z\) is a
probability law, as \(a_z z_P\) are mutually orthogonal projections
which sum to \(a_z\). Haar normalization divides all identity
coefficients of the finite group algebra by the measure of the
principal congruence kernel. It cancels in the ratio, giving
\[
 M_i(S)=\frac{e_{i,S}(1)}{e(1)}
       =\sum_{\boldsymbol{P}\in S}\prod_{z\mid g_i}\omega_z(P_z).
\]
Each \(e_{i,S}\) has the same principal congruence level as \(e\) in the
sense of bi-invariance. It is self-adjoint, idempotent and of positive
type. For example, a self-adjoint convolution idempotent \(a\) of
compact support satisfies \(a(1)=\|a\|_2^2\), and the kernel
\(a(x^{-1}y)\) is the Gram kernel of its right translates in \(L^2\).
The finite support of \(e_{i,S}\) is integral and is unipotent modulo
\(g_{1-i}\), since the factor at this prime is unaltered. Thus the
\hyperlink{candidate.count}{multiplicity estimate} applies with
\(g_{1-i}\) as the support prime and the same real data as for \(e\).
If \(D_{i,S}\) is its multiplicity and
\(\eta_c=e^{-rf\log r/16}\), then
\[
  D_{i,S}\le (1+\eta_c)D_\star M_i(S),\qquad
  D\ge(1-\eta_c)D_\star .
\]
This remains valid when \(e_{i,S}=0\).

\paragraph{The local character projectors.}
Write \(s_0=m/2-1\). The actual local factor \(a_z\) is the
character projector of the upper unitriangular subgroup \(J\)
defined by the following block order of lengths \(s_0,k,1,s_0\).
Call the middle \(k\) indices \(B\), and let \(t\) be the singleton
immediately following \(B\). Entries from \(B\) to column \(t\) are
zero. The additive character \(\psi_g\ne1\) of \(\mathbf F_g\)
is applied to the sum of consecutive entries on the \(B\)-chain and
on the complementary chain, which passes through \(t\).
The group has index \(g^k\) in the upper unitriangular group \(U_r\).
These are \(e_J\) and its ordinary Whittaker analogue \(e_U\) in the
\releasecite{Formalism/SharpRestriction/LevelLetterCocycles/ArithmeticCylinders.lean}{298}{finite model}.
In the adapted split bases the row level uses \(e_J\) at every
place over the two auxiliary primes. A change of basis conjugates
\(e_J\) and leaves its mass against the central \(z_P\) unchanged.
The argument applies to every nontrivial residue-field additive
character, including those obtained in the different local bases.
\releasecite{Formalism/LocalToGlobal/SiegelLatticeTheory/FlagStabilizerTensors/RootLatticeThresholds/AllRowsKpp.lean}{22}{Adapted split bases}
\releasecite{Formalism/LocalToGlobal/SiegelLatticeTheory/FlagStabilizerTensors/RootLatticeThresholds/AllRowsEtOne.lean}{104}{Identity coefficients}
\releasecite{Formalism/LocalToGlobal/SiegelLatticeTheory/FlagStabilizerTensors/AllRowsOf.lean}{512}{Projectors at all places}
\releasecite{Development/AbelianVarietyTheory/Proofs/FlagMeasures/AllRowsAtRecord.lean}{67}{Their common row record}.
The ordinary Whittaker law \(\omega_U(P)=(e_U z_P)[1]/e_U[1]\)
is a probability law by the same trace argument. We show that
\begin{equation}
\label{eq:trace-finite-mass-comparison}
  \omega_z(P)\ \ge\ (1-1/g)^k\omega_U(P)\qquad\text{for every }P .
\end{equation}

To see the scalar in \eqref{eq:trace-finite-mass-comparison}, first
move the left-hand \(s_0\) indices past \(B\), one at a time. In one
step let \(a\) be the last index before \(B\), and \(t\) the first of
the remaining indices after it. The missing column from \(B\) is at
\(t\), and the complementary character uses the entry \(a\to t\).
Take \(X\) to be the root groups in the row \(a\to B\), and \(Y\)
those in the column \(B\to t\). Remove \(X\) and the internal
\(U_B\) from the current pattern to obtain the common subgroup \(N\).
The subgroups \(X,Y,U_B\) normalize \(N\) and preserve its character.
In matrix multiplication the changes in character coordinates can
only involve \(a\to t\). Conjugation by \(X\) cannot change this entry
because \(B\to t\) is zero in \(N\). Conjugation by \(Y\) cannot
change it because \(a\to B\) is zero there. Conjugation by \(U_B\)
fixes all entries outside the rows and columns of \(B\).
For a representation \(V\), on \(V^{N,\psi_g}\) the commutator of
x in \(X\) and y in \(Y\) acts by the nondegenerate pairing
\(\psi_g(xy)\) of a row and a column. The averaging operators \(P_X,P_Y\)
satisfy \(P_XP_Y=|Y|^{-1}I\) on \(X\)-invariants and
\(P_YP_X=|X|^{-1}I\) on \(Y\)-invariants, by character orthogonality.
They commute with \(U_B\), which normalizes \(X\) and \(Y\).
Hence replacing \(X\) by \(Y\) preserves the character multiplicity
including the Whittaker condition in \(U_B\). In the new order
with \(a\) after \(B\), this makes \(B\to a\) the missing column.
Repeating the exchange leaves \(B\) first and the complementary
chain as a single tail.

For that final order, first take invariants of the cross-block
columns other than the first tail column, then the tail's Whittaker
eigenspace. The resulting space \(V_0\) carries the group
\(P_k=\mathrm{GL}_k(\mathbf F_g)\ltimes\mathbf F_g^k\), where the
translation group is the first cross-block column. Indeed the other
cross columns are normal in the relevant block-upper subgroup and
the tail unipotent group centralizes the first cross-block column
modulo those columns. The rank of the new \(J\)-projector is the
dimension of the \(U_k\)-Whittaker space in \(V_0\). The full
\(U_r\)-Whittaker rank is the dimension with the additional
translation character \(z\mapsto\psi_g(z_k)\).
By Frobenius reciprocity this latter rank counts copies in \(V_0\)
of the irreducible mirabolic representation
\[
 T_0=\operatorname{Ind}_{\,U_k\ltimes\mathbf F_g^k}^{\,P_k}
       (\psi_{U_k}\otimes\psi_g(z_k)).
\]
\cite[Proposition~1.2, p.~18]{tr3gelfand}
The restriction of \(T_0\) to \(\mathrm{GL}_k\) is
\(\operatorname{Ind}_{U_k}^{\mathrm{GL}_k}\psi_{U_k}\), as follows by
using \(\mathrm{GL}_k\) as representatives modulo translations in
the induction model. Its \(U_k\)-Whittaker dimension is at least
\((g-1)^k\). For \(d\) diagonal and \(w_0\) the longest permutation,
the double cosets \(U_k d w_0 U_k\) have trivial stabilizer
intersection because \(w_0 U_k w_0^{-1}\) is lower unitriangular.
Distinct \(d,d'\) give distinct cosets, since an upper triangular
matrix with diagonal \(d'^{-1}d\) can be lower unitriangular only
when \(d=d'\). Each coset thus supports a Whittaker function.
It follows for every \(V\) that
\(\operatorname{rank} e_J|_V\ge(g-1)^k\operatorname{rank} e_U|_V\).
Apply this to the regular representation on the \(z_P\)-part and use
\(e_J[1]/e_U[1]=g^k\) to obtain \eqref{eq:trace-finite-mass-comparison}.

\paragraph{The squarefree law.}
On a squarefree label \(P=\prod_h P_h\), with distinct irreducible
factors of degrees \(a_h\), Green's character formula gives a unique
irreducible in the label and degree
\(\prod_{j=1}^r(g^j-1)/\prod_h(g^{a_h}-1)\).
It is the parabolic induction of the cuspidal representations of
the distinct factors. The formulas apply with the label convention
of the period spaces, since renaming cuspidal parameters within
each degree preserves the squarefree formula.
\cite[Introduction, \S7 Theorem~12 and \S8 Theorems~13--14]{tr3green}
Each cuspidal factor has a Whittaker line.
\cite[Theorem~0.5 and Theorems~2.3 and~3.3]{tr3gelfand}
The open parabolic--unipotent double coset in the induced model
supplies the product of those lines, and the full Whittaker rank
is at most one by Gel'fand's Theorem~0.5. Applying the regular
representation trace calculation to this irreducible, and using
\([\mathrm{GL}_r(\mathbf F_g):U_r]=\prod_{j=1}^r(g^j-1)\), gives
\[
 \omega_U(P)=\prod_h (g^{a_h}-1)^{-1}.
\]

There is a direct way to use these exact masses. For one place,
draw a uniform permutation of \(r\) indices and, independently for
each cycle of length \(a\), draw a uniform element of
\(\mathbf F_{g^a}^{\times}\). Multiply the monic linear factors for
all its \(a\) Frobenius conjugates, including repetitions when its
degree is smaller than \(a\). For prescribed squarefree \(P\) with
\(n_a\) factors of degree \(a\), the permutation type has probability
\(\prod_a (a^{n_a}n_a!)^{-1}\). Assign the specified factors to its
cycles in \(\prod_a n_a!\) ways and choose one of the \(a\) roots for
each cycle. The resulting probability of \(P\) is exactly
\(\prod_h(g^{a_h}-1)^{-1}\). Use independent draws at the \(f\) places.
The product laws have the same probability for every tuple that is
squarefree at all places, so they give the same probability to any
condition implying squarefreeness at all places.

\paragraph{The bad-label mass.}
At each \(g=g_i\) the good tuples are squarefree at all places,
their distinct roots stay distinct after taking their prime-to-\(\ell\)
parts for all primes \(\ell\notin\mathcal B(r,g)\), and no root
agrees in prime-to-\(\ell\) part with any of the \(m\) prescribed
twisting values at its place or their inverses for any \(\ell\ne g\).
Here
\[
 \mathcal B(r,g)=\{g\}\cup
   \{\ell:\ell\text{ prime},\ \ell^{4v_\ell(g^a-1)}>g^a
                       \text{ for some }1\le a\le r\}.
\]
The tuple must also have no common divisor degree in
\([1,\lfloor r/2\rfloor]\), and have a factor of prime degree \(p\)
with \(r/2<p\le r-3\) at some place.
\releasecite{Formalism/CyclotomicNorms/StatementT.lean}{19}{Good label conditions}
and \releasecite{Formalism/SharpRestriction/LevelLetterCocycles/ArithmeticCylinders.lean}{355}{both twist signs}.

In the permutation model the probability of a repeated root at a
place is \(O(r^3/g)\). A cycle of length \(a\ge2\) lands in a proper
subfield with probability at most \(2a g^{-a/2}\). When all cycles
draw elements of their full degrees, a repeat between two cycles
requires equal lengths \(a\), and then there are \(a\) choices for
the second root out of \(g^a-1\). These bounds hold conditionally on
the entire permutation. For distinct values \(u,v\), equality
of prime-to-\(\ell\) parts means \(u/v\) has nontrivial \(\ell\)-power
order. For roots from two cycles of lengths \(a,a'\), the prime
must divide \((g^a-1)(g^{a'}-1)\). With the first value fixed, the
possible second values occupy at most one coset of the \(\ell\)-primary
subgroup of \(\mathbf F_{g^{a'}}^\times\). For \(\ell\notin\mathcal B(r,g)\)
its size is at most \(g^{a'/4}\). For roots in the same cycle of
length \(a\), their ratio has the form \(u^{g^j-1}\), \(0<j<a\).
There are at most \(g^{a/2}g^{a/4}\) choices for \(u\) with this ratio
of \(\ell\)-power order, since the kernel of this power map has order
\(g^{\gcd(a,j)}-1\). Dividing by \(g^a-1\) and summing over at
most \(2r\log_2 g\) eligible primes and \(r^2\) pairs gives
\(O(r^3\log g/\sqrt g)\) for these collisions.

Each specified twist value is in \(\mathbf F_g^\times\). The equality
test against a root from a cycle of length \(a\) has probability
at most \((g^a-1)_\ell/(g^a-1)\), with \((n)_\ell\) the
\(\ell\)-primary factor. Since \(q,d\) are distinct primes and
\(qd\mid g-1\), at least one of them divides the complementary factor.
Thus the probability is at most \(1/\min(q,d)\) for each
\(\ell\mid g^a-1\). It suffices to consider these primes, because
equality of the root and the twist value is already detected by any
one of them. A union bound over roots, primes and the \(2m\) values
gives \(O(r^2m\log g/\min(q,d))\) at a place.

On squarefree outcomes, divisor degrees are the sizes of unions
of cycles. For \(1\le a\le r/2\) let \(Z_a\) be the number of invariant
\(a\)-subsets in a uniform permutation. Its expectation is \(1\).
Size-biasing by \(Z_a\) is achieved by prescribing a uniform \(a\)-subset
and drawing independent uniform permutations on it and its complement.
A uniform permutation on a nonempty set has a fixed point with
probability at least \(1/2\), by the first two inclusion--exclusion
terms for the events that individual points are fixed. With probability
at least \(1/4\) in the size-biased law there are fixed points on both
sides, and swapping them gives a second invariant \(a\)-subset.
Consequently \(\Pr(Z_a=0)=\mathbb E (Z_a-1)_+\ge
\tfrac12\mathbb E[Z_a\mathbf 1_{Z_a\ge2}]\ge1/8\).
Independence gives at most \(r(7/8)^f\) for the obstruction of any degree.
A prime cycle \(r/2<p\le r-3\) has probability \(1/p\),
and these events are disjoint for different \(p\).
Their sum is at least \(c/\log r\) for an absolute \(c>0\) and all
large \(r\), by the lower bound for \(\pi(r)-\pi(r/2)\) with at
most three primes removed.
\cite[Corollary~3, (3.8), p.~69]{rosser-schoenfeld}
Failure at all places therefore costs at most \(\exp(-cf/\log r)\).

For the actual product laws the comparison
\eqref{eq:trace-finite-mass-comparison} gives at least
\((1-1/g)^{kf}\) times the ordinary product-law mass of the good set.
Let \(S_i^{\rm bad}\) be the complementary label set. Bernoulli's
inequality and the preceding bounds give
\begin{equation}
\label{eq:trace-bad-mass}
 \begin{split}
 M_i(S_i^{\rm bad})\ \ll\
 &\frac{kf}{g_i}
 +\frac{fr^3\log g_i}{\sqrt{g_i}}
 +\frac{fr^2m\log g_i}{\min(q,d)}\\
 &+r(7/8)^f+\exp(-cf/\log r).
 \end{split}
\end{equation}
The comparison uses both internal Whittaker conditions of the
actual finite projector.

The selected sizes give \(f\ge x^2\), \(m,f,r\le x^6\),
\(g_i\ge\exp(8x^{\varepsilon_{\rm row}})\),
\(\log g_i\ll x^\varepsilon\), and
\(\min(q,d)\ge\exp(x^{\varepsilon_{\rm row}/6})\), where
\(\varepsilon_{\rm row}=\varepsilon/J>0\) is fixed.
Thus each \(M_i(S_i^{\rm bad})\) is at most a polynomial in \(x\)
times \(\exp(-x^{\varepsilon_{\rm row}/6})\). The complement of
the common good space lies in the sum of the two bad ranges. Counting
both ranges by their dimensions, including all repetitions, gives
\[
 \frac{D-D'}D\le
 \frac{1+\eta_c}{1-\eta_c}\bigl(M_0(S_0^{\rm bad})+M_1(S_1^{\rm bad})\bigr).
\]
Since \(\eta_c\to0\) uniformly and
\(x^{\varepsilon_{\rm row}/6}-x^{\varepsilon_{\rm row}/7}\) dominates
every fixed multiple of \(\log x\), this proves
\begin{equation}
 \label{eq:trace-row-fraction}
 D'=(1-\delta)D,\qquad
 0\le\delta\le\exp(-x^{\varepsilon_{\rm row}/7})\le\tfrac12 .
\end{equation}

\subsection{From ranks to the mean slope}
\leavevmode\teamanchor{proof.height-filtration}%
For each word use \(N'_w=I_wN_w\) and \(C=R^\vee\) from
\eqref{eq:period-lattice}. Multiplication by a nonzero fractional
ideal \(I_w\) preserves the \(K\)-span of \(N_w\). Thus \(N'_w\)
has rank \(r_w\) and its \(K\)-space has the prescribed realization
\(\mathcal M_w\) at \(\iota_K\). Division in Section~\ref{sec:integral}
gives \(B_{K,w}(N'_w)\subseteq C\). The comparison below uses
the squared-norm archimedean height bound \(h_w\) for this same
map in the prescribed metrics on \(N'_w\) and \(C\).

Use the squared-norm slope
\(\mu(L)=2\widehat{\deg}_{\rm n}(L)/\operatorname{rk}_{\mathcal O_K}L\),
where \(\widehat{\deg}_{\rm n}\) is the Arakelov degree divided by
\([K:\mathbf Q]\). In the Euclidean realization with one copy of
\(\mathbf C^{\operatorname{rk}L}\) for each complex place, the
release doubles the slope from \(-\log\) covolume per real dimension
and adds \(-\log 2+\log|\operatorname{disc}K|/[K:\mathbf Q]\).
Indeed, the reference covolume per \(K\)-coordinate is
\(2^{-[K:\mathbf Q]/2}\sqrt{|\operatorname{disc}K|}\).
\releasecite{Formalism/HarderNarasimhan/HarderNarasimhanProofs/FrameHermSharpB.lean}{80}{conversion from Euclidean covolume to the normalized slope}
Equip the \(K\)-spans of the Hermitian lattices
with their given archimedean norms and their lattice unit balls at the
finite places. For a nonzero \(K\)-linear map \(B\) between these
Hermitian adelic bundles, write \(B_\nu\) for its scalar extension to
the completion \(K_\nu\). Its normalized adelic and squared-norm
archimedean heights are
\[
  h_{\rm ad}(B)=\frac{1}{[K:\mathbf Q]}
       \sum_\nu n_\nu\log\|B_\nu\|_{\mathrm{op}},
  \qquad
  h_\infty(B)=\frac{1}{n_\infty}
       \sum_{\nu\mid\infty}\log\bigl(\|B_\nu\|_{\mathrm{op}}^2\bigr).
\]
Here \(n_\nu=[K_\nu:\mathbf Q_p]\) for \(\nu\mid p\) and
\(n_\nu=2\) for the \(n_\infty=[K:\mathbf Q]/2\) complex places.
The absolute values at finite places extend the usual \(p\)-adic ones.
\cite[Definition~6.1, arXiv version]{rk3gaudron}
\cite[Equation~(20), p.~34]{tr3chenHN}
Rescaling an arithmetic map by \(\lambda\in K^\times\) multiplies its
image lattice by \(\lambda\) and adds
\(2\log|\mathrm{N}_{K/\mathbf Q}\lambda|/[K:\mathbf Q]\)
to \(h_\infty(B)\). For a unit of \(\mathcal O_K\) both the lattice image
and the archimedean height are preserved. The product formula gives
\(h_{\rm ad}(\lambda B)=h_{\rm ad}(B)\) for every \(\lambda\in K^\times\).

For a Hermitian lattice \(L\) and any saturated sublattice \(M\) with
the induced norms, \(\widehat{\deg}_{\rm n}(L)=
\widehat{\deg}_{\rm n}(M)+\widehat{\deg}_{\rm n}(L/M)\) with the
quotient norms on \(L/M\).
The Harder--Narasimhan polygon of \(L\) is the upper
concave hull of \((\operatorname{rk}M,2\widehat{\deg}_{\rm n}(M))\)
over these sublattices, including \(0\) and \(L\).
There is a unique saturated flag whose rank-degree pairs are exactly
the vertices. The slopes of the successive quotient lattices are the
segment slopes in decreasing order.
\cite[Propositions~4.22 and~5.13, arXiv version]{rk3gaudron}
Write \(\lambda_j(L)\) for these slopes, repeated according to the
ranks of the quotients. Their mean is \(\mu(L)\) and their extrema
are \(\mu_{\min}(L)\) and \(\mu_{\max}(L)\).
The saturated step \(L_{\ge a}\) retains the factors of slope at least
\(a\) and has rank \(\#\{j:\lambda_j(L)\ge a\}\).

\leavevmode\teamanchor{proof.determinant-height-comparison}%
For the usual slopes, which are half the squared-norm slopes, put
\(\mathcal F^L_u=L_{\ge 2u}\otimes_{\mathcal O_K}K\) in
\(L_K=L\otimes_{\mathcal O_K}K\) and define
\(\mathcal F^C_u\) in the same way. Chen's filtration shift gives,
for every nonzero \(K\)-linear map \(B:L_K\to C_K\) and every
\(u\in\mathbf R\),
\[
 B(\mathcal F^L_u)\subseteq \mathcal F^C_{u-h_{\rm ad}(B)}.
\]
\cite[Proposition~2.2.4, p.~42]{tr3chenHN}
If \(B(L)\subseteq C\), the finite operator norms are at most \(1\),
so \(2h_{\rm ad}(B)\le h_\infty(B)\). If \(h_\infty(B)\le h\), apply
the inclusion at \(u=(a+h)/2\) and intersect with \(C\) to obtain
\(B(L_{\ge a+h})\subseteq C_{\ge a}\) for every \(a\in\mathbf R\).
The lattice inclusion also holds for the zero map.

\leavevmode\teamanchor{proof.integrated-slope-deficit}%
Choose uniformly from the \(R_{\rm act}\) donor slopes, with repetitions
according to rank. The adjusted value is \(X=\lambda_j(N'_w)-h_w\)
when the chosen slope is the \(j\)-th of \(N'_w\). Set
\[
 s=\mathbb E X=\frac{\sum_w r_w(\mu(N'_w)-h_w)}{R_{\rm act}},
 \qquad U=\max X-s\ge0 .
\]
For \(\tau\ge0\), take the steps \((N'_w)_{\ge s-\tau+h_w}\).
The discarded rank fraction is \(\alpha_\tau=\Pr(X<s-\tau)\).
The retained periods land in \(C_{\ge s-\tau}\). Their \(K\)-spans
are subspaces of the same actual column models.
By \eqref{eq:trace-rational-corank}, the fraction of target slopes
below \(s-\tau\), denoted by \(\beta_\tau\), is at most
\((\eta+4\alpha_\tau+4\sqrt e)/(1-\delta)\).
Choose \(H_*\ge (s-\mu_{\min}(C))_+\). Integrating the two tails gives
\[
 \begin{aligned}
 s-\mu(C)&\le \frac1{D'}\sum_{j=1}^{D'}(s-\lambda_j(C))_+
   =\int_0^{H_*}\beta_\tau\,d\tau,\\
 \int_0^{H_*}\alpha_\tau\,d\tau&=\mathbb E\min\{H_*,(s-X)_+\}.
 \end{aligned}
\]
Write \(Z=s+U-X\ge0\), with \(\mathbb E Z=U\). If \(H_*+U>0\),
\[
 \min\{H_*,(Z-U)_+\}\le \frac{H_* Z}{U+H_*}.
\]
For \(Z\le U+H_*\), the right side bounds both \(0\) and \(Z-U\).
For \(Z\ge U+H_*\) it bounds \(H_*\). Thus
\[
 s-\mu(C)\le
 \frac{4UH_*/(U+H_*)+(\eta+4\sqrt e)H_*}{1-\delta}.
\]
The fraction \(UH_*/(U+H_*)\) is zero when either variable is zero,
including their common zero. At \(U=0\) every adjusted donor slope
equals \(s\). At \(H_*=0\) every target slope is at least \(s\).

In particular \(UH_*/(U+H_*)\le U\). For \(U>0\), one also has
\(UH_*/(U+H_*)\le U\log(1+H_*/U)\), by integrating
\(U/(U+\tau)\) on \([0,H_*]\). The
\releasecite{Formalism/HarderNarasimhan/HarderNarasimhanProofs/TheoremFSharp.lean}{390}{form used in the arithmetic comparison}
uses this logarithmic bound and adds the nonnegative \(\log|W|\).

\begin{equation}
\teamanchor{proof.conditional-slope-bound}
\label{eq:trace-slope}
 s-\mu(C)\ \le\ \log|W|+
 \frac{4U\log(1+H_*/U)+(\eta+4\sqrt e)H_*}{1-\delta}.
\end{equation}
The \(U\)-term is zero at \(U=0\). The slopes of a Hermitian dual
are the negatives of the original slopes in reverse order.
\cite[Properties~5.14(2), Hermitian case, arXiv version]{rk3gaudron}
This gives \(s-\mu_{\min}(C)=s+\mu_{\max}(R)\) in
Section~\ref{sec:integral}.
The height comparison there bounds this quantity by \(x^{C_H}b\)
with the same \(C_H\) as in the word count. Taking
\(H_*=\max(0,s-\mu_{\min}(C))\), the numerator satisfies
\((\eta+4\sqrt e)H_*\le6b\) by \eqref{eq:trace-accuracy}.
Since \(1-\delta\ge1/2\) by \eqref{eq:trace-row-fraction}, this term
contributes at most \(12b\). Using \(UH_*/(U+H_*)\le U\) also gives
\begin{equation}
\label{eq:trace-sharp-slope}
  s-\mu(C)\le 8U+12b.
\end{equation}
